% paper_v2_vn.tex -- the Vietnamese version of paper_v2: the same paper in ACM's two-column format
% (acmart, sigconf), translated in parts L1-L9 (plan_v2_vn.tex in this folder). ../paper_v2/ is the
% source and is not edited for this one (N1); this folder has its own sections/, tables/, figures/
% and numbers.tex, with the same file names as paper_v2's.
% Since N31 the paper has no appendix: the former appendix is the supplementary document,
% supplement_vn.tex, which is built from the same folder after this file (latexmkrc builds both).
% Compile with LuaLaTeX and BibTeX (in Overleaf: Menu, Compiler, LuaLaTeX): under pdfLaTeX,
% acmart's fonts have no Vietnamese letters. It uses the three .bib files at the project root.
% numbers.tex and figures/*.dat: from the LMSPS repository's root, with the HAN repository next
% to it, python paper/make_numbers.py --lmsps . --han ../han-improve --out <this folder> --lang vi
\RequirePackage{iftex}
\ifLuaTeX\else
  \errmessage{paper_v2_vn.tex needs LuaLaTeX. In Overleaf: Menu, Compiler, LuaLaTeX}%
  \expandafter\stop
\fi
\PassOptionsToPackage{quiet}{fontspec}   % notes on font features go to the log, not the warnings
\PassOptionsToPackage{dvipsnames}{xcolor}
\PassOptionsToPackage{obeyspaces,spaces}{url}   % \code keeps its spaces and can break at them
% nonacm: a draft not yet at an ACM venue, so no conference line, copyright block or ACM reference
% box.
\documentclass[sigconf,nonacm]{acmart}
\makeatletter
\@ACM@balancefalse   % no column balancing: acmart's \balance would come in the second column
\makeatother         % of the last page and only warn that it cannot balance there
\raggedbottom

% Vietnamese is the main language: babel's captions (Hình, Bảng, Phụ lục) and line breaking.
% English stays loaded for the reference list, which keeps the language of the works (N7), for
% ACM's CCS terms, and for the English words of quotations (\foreignlanguage{english}, N13).
\usepackage[english,vietnamese]{babel}
\providecommand\ccsname{CCS Concepts}
\addto\captionsvietnamese{%
  \renewcommand\abstractname{Tóm tắt}%          babel's is Tóm tắt nội dung
  \renewcommand\keywordsname{Từ khóa}%
  \renewcommand\acksname{Lời cảm ơn}%
  \renewcommand\refname{Tài liệu tham khảo}%    babel's is only Tài liệu
  \renewcommand\ccsname{Phân loại CCS}}
% acmart prints the heading of the CCS concepts as fixed text; this makes it \ccsname, and sets
% ACM's terms under it as English.
\usepackage{etoolbox}
\makeatletter
\patchcmd{\maketitle}{\@specialsection{CCS Concepts}%
         \noindent\@concepts\par}{\@specialsection{\ccsname}%
         \noindent\foreignlanguage{english}{\@concepts}\par}{}%
  {\PackageWarningNoLine{paper_v2_vn}{acmart's CCS block has changed, so its heading stays
     English; see \string\ccsname}}   % the same code in acmart 2.03 (2024) and 2.21 (2026)
\makeatother
\setlength{\emergencystretch}{1em}   % Vietnamese words are not hyphenated: a little more stretch
\AtEndPreamble{\hypersetup{pdflang={vi}}}   % the PDF's language; runs after acmart's pdflang={en}

\usepackage{array}
\usepackage{tabularx}
\usepackage{tikz}
\usetikzlibrary{arrows.meta,positioning}
\usepackage{pgfplots}
\pgfplotsset{compat=1.18}
\pgfkeys{/pgf/number format/use comma}   % axis numbers as 0,05 and 1.000, as in numbers.tex
\usepackage{icomma}   % in a formula, a comma between digits is a decimal comma: $-0,034$

\usepackage{xurl}   % long links may break anywhere in the narrow columns
% File names and meta-paths in typewriter type; they may break only after / _ . or at a space, so
% a meta-path such as PUPUP is never split.
\DeclareUrlCommand\codeurl{\def\UrlBreaks{\do\/\do\_\do\.}\def\UrlBigBreaks{}\urlstyle{tt}}
\DeclareRobustCommand{\code}[1]{\codeurl{#1}}
\newcolumntype{L}{>{\raggedright\arraybackslash}X}
% Wide tables and figures span both columns.
\newenvironment{widetable}[1][tbp]{\begin{table*}[#1]}{\end{table*}}
\newenvironment{widefigure}[1][tbp]{\begin{figure*}[#1]}{\end{figure*}}
% Figure colours, as in paper_v2: searched sets blue, manual orange, random draws aqua.
\definecolor{cSearch}{HTML}{2A78D6}
\definecolor{cManual}{HTML}{EB6834}
\definecolor{cRandom}{HTML}{1BAF7A}
\definecolor{cOther}{HTML}{7C7B76}
% Every number of the paper: \val{key}, from numbers.tex (Vietnamese number format, N4).
% numbers.tex -- written by make_numbers.py; do not edit by hand, run the script again.
% 2026-11-10: python make_numbers.py --lmsps LMSPS --han han-improve --lang vi
% The raw-file counts come from extraction_report.json; the block below names the copy used.
% Vietnamese number format (--lang vi): 0,666 and 1.147.920; versions, commits and path sets as in
% English. The notes after the values stay in English.
% \val{key} prints a number; an unknown key prints ??key in bold and raises a LaTeX warning.
\makeatletter
\providecommand{\setval}[2]{\@namedef{val@#1}{#2}}
\DeclareRobustCommand{\val}[1]{\@ifundefined{val@#1}{\textbf{??#1}%
  \PackageWarning{numbers}{No value `#1' in numbers.tex}}{\@nameuse{val@#1}}}
\makeatother

% --- Data: the raw file, the fixes and the clean graph (LMSPS notes, Parts 2-4) ------------------
\setval{raw-records}{548.552}  % extraction_report.json (review/evidence/extraction_report.json)
\setval{raw-usable}{542.684}  % extraction_report.json (review/evidence/extraction_report.json)
\setval{raw-reviews}{7.593.244}  % extraction_report.json (review/evidence/extraction_report.json)
\setval{raw-categories}{49.732}  % extraction_report.json (review/evidence/extraction_report.json)
\setval{raw-pos-pct}{78,6}  % audit_before_fix.json: share of review lines rated 4 or 5
\setval{raw-unhelpful-pct}{1,35}  % positive reviews dropped as unhelpful (>= 5 votes, < 20% helpful)
\setval{placeholder-ids}{6}  % extraction_report.json (review/evidence/extraction_report.json)
\setval{placeholder-lines}{1.147.920}  % extraction_report.json (review/evidence/extraction_report.json)
\setval{placeholder-main-lines}{945.065}  % extraction_report.json (review/evidence/extraction_report.json): ATVPDKIKX0DER
\setval{placeholder-second-lines}{201.770}  % extraction_report.json (review/evidence/extraction_report.json): A3UN6WX5RRO2AG
\setval{edition-pairs}{64.814}  % extraction_report.json (review/evidence/extraction_report.json)
\setval{edition-clusters}{30.648}  % extraction_report.json (review/evidence/extraction_report.json)
\setval{edition-products}{72.280}  % extraction_report.json (review/evidence/extraction_report.json)
\setval{raw-works}{501.052}  % usable products after merging editions
\setval{kcore-iterations}{13}  % extraction_report.json (review/evidence/extraction_report.json)
\setval{catpath-rows}{2.288.922}  % extraction_report.json (review/evidence/extraction_report.json)
\setval{placeholder-main-max}{1.056}  % most reviews on one product
\setval{placeholder-second-max}{1.193}
\setval{placeholder-found}{4}  % IDs found by the rule, not named in advance
\setval{placeholder-found-max-lo}{57}  % their most reviews on one product
\setval{placeholder-found-max-hi}{149}
\setval{placeholder-found-lines}{1.085}
\setval{edition-size-min}{2}  % products in an edition cluster
\setval{edition-size-max}{24}
\setval{edition-size-two-pct}{79,5}  % clusters of exactly two products
\setval{rule-pos-rating}{4}  % a positive review is rated at least this
\setval{rule-helpful-votes}{5}  % a review with at least this many votes ...
\setval{rule-helpful-pct}{20}  % ... and fewer helpful ones than this is dropped
\setval{rule-kcore}{5}
\setval{rule-edition-shared}{3}  % identical review lines two editions must share ...
\setval{rule-edition-overlap-pct}{50}  % ... covering this share of the smaller one
\setval{rule-edition-maxrun}{50}  % review lines on more products than this are ignored
\setval{rule-placeholder-repeats}{50}  % reviews on one product that mark a placeholder
\setval{before-products}{90.659}
\setval{before-categories}{39.754}
\setval{before-users}{171.285}
\setval{before-edges}{2.612.766}
\setval{before-maxdeg}{55.572}  % largest user degree
\setval{before-meddeg}{7}  % median user degree
\setval{before-clusters}{17.726}  % duplicate-edition clusters
\setval{before-clustered}{44.591}  % products in those clusters
\setval{before-clustered-pct}{49,2}
\setval{before-copyedges-pct}{55,0}  % edges that copy an interaction already in the graph
\setval{before-valleak-pct}{72,9}  % validation edges with a training copy
\setval{before-testleak-pct}{73,0}  % test edges with a training copy
\setval{before-valtitle-pct}{60,3}  % validation edges whose user has a training edge to a same-title product
\setval{before-testtitle-pct}{60,2}  % the same for test edges
\setval{before-alsobought}{188.257}  % distinct also-bought product pairs
\setval{after-products}{48.883}
\setval{after-categories}{37.411}
\setval{after-users}{53.569}
\setval{after-edges}{819.847}
\setval{after-maxdeg}{2.134}  % largest user degree
\setval{after-meddeg}{8}  % median user degree
\setval{after-clusters}{6}  % duplicate-edition clusters
\setval{after-clustered}{12}  % products in those clusters
\setval{after-clustered-pct}{0,025}
\setval{after-copyedges-pct}{0,007}  % edges that copy an interaction already in the graph
\setval{after-valleak-pct}{0,011}  % validation edges with a training copy
\setval{after-testleak-pct}{0,010}  % test edges with a training copy
\setval{after-valtitle-pct}{0,86}  % validation edges whose user has a training edge to a same-title product
\setval{after-testtitle-pct}{0,80}  % the same for test edges
\setval{after-alsobought}{87.901}  % distinct also-bought product pairs
\setval{after-zerocat-pct}{60,7}  % categories with an all-zero name vector
\setval{after-leafcats-med}{5}  % leaf categories per product
\setval{title-coverage-pct}{95,8}  % products with a non-zero title vector
\setval{title-dim}{256}
\setval{title-vocab}{12.699}
\setval{cat-vocab}{128}
\setval{user-bits}{8}
\setval{user-dim}{64}
\setval{rel-pp-rows}{121.075}  % also-bought rows per direction
\setval{rel-cc}{37.400}
\setval{rel-pc}{427.565}
\setval{rel-up}{819.847}

% --- The split (rebuilt as hgb/link_data.py builds it; checked against the fingerprint) ----------
\setval{split-train}{655.879}
\setval{split-val}{81.984}
\setval{split-test}{81.984}
\setval{split-seed}{1}
\setval{split-fingerprint}{c5f7bccbf61c1b0f}
\setval{split-negatives}{99}
\setval{train-deg-med}{6}  % training edges per user
\setval{train-deg-mean}{12,2}
\setval{train-deg-max}{1.745}
\setval{test-users}{33.222}  % users with at least one test edge
\setval{group-book}{23.754}
\setval{group-book-pct}{48,6}
\setval{test-book-pct}{41,0}  % share of test edges
\setval{group-dvd}{4.359}
\setval{group-dvd-pct}{8,9}
\setval{test-dvd-pct}{13,9}  % share of test edges
\setval{group-music}{16.436}
\setval{group-music-pct}{33,6}
\setval{test-music-pct}{33,8}  % share of test edges
\setval{group-video}{4.334}
\setval{group-video-pct}{8,9}
\setval{test-video-pct}{11,3}  % share of test edges
\setval{group-bookmusic-pct}{82,2}
\setval{test-bookmusic-pct}{74,8}  % books and music, share of test edges

% --- The graph: category links, near-empty paths and edition twins (audits, link.dat, HAN logs) ---
\setval{before-topuser-pct}{61}  % share of products the largest user had reviewed before the fix
\setval{raw-shared-pct}{68,7}  % raw review lines that also appear under another product
\setval{cat-all-median}{13}  % categories per product if every ancestor were linked
\setval{cat-all-largest-pct}{48,7}  % products in one category, every ancestor linked
\setval{cat-pcp-ratio}{20,7}  % P-C-P walks with every ancestor linked, over leaves only
\setval{pcp-isolated}{9}  % products whose leaf categories hold no other product
\setval{pccp-reach-train-pct}{27,7}  % training edges whose product PCCP links to another product
\setval{pccp-isolated-pct}{81,1}  % products PCCP links to no other product (full graph)
\setval{pccp-reach-book-pct}{2,2}  % share of Book products that PCCP links to another product
\setval{pccp-reach-dvd-pct}{90,1}  % share of DVD products that PCCP links to another product
\setval{pccp-reach-music-pct}{5,3}  % share of Music products that PCCP links to another product
\setval{pccp-reach-video-pct}{90,2}  % share of Video products that PCCP links to another product
\setval{empty-isolated-min-pct}{76,2}  % near-empty paths: products without a neighbour, HAN graphs
\setval{empty-isolated-max-pct}{81,1}
\setval{pccp-graph-isolated-pct}{81,1}  % PCCP in the HAN graph (pruning adds a few)
\setval{pcccp-graph-isolated}{0}
\setval{pcp-graph-isolated}{9}
\setval{pp-isolated-pct}{11,3}  % products without an also-bought partner
\setval{twin-products}{51}  % products with an identical reviewer set elsewhere
\setval{twin-groups}{23}
\setval{twin-edges}{497}
\setval{twin-val}{47}
\setval{twin-test}{40}
\setval{example-shared}{659}  % review lines a DVD and a VHS of one film share (review part 1)
\setval{example-lines}{660}

% --- The search: 30 runs per set size (LMSPS notes, Parts 1 and 5-9) -----------------------------
\setval{search-cands}{36}
\setval{search-pool}{10}
\setval{search-sets}{210}  % sets scored per run (10 choose 4 = 10 choose 6)
\setval{search-epochs}{150}
\setval{search-warmup}{15}
\setval{search-hops}{5}
\setval{search-hidden}{512}
\setval{search-batch}{10.000}
\setval{search-dropout}{0,5}
\setval{search-lr}{0,001}
\setval{search-tau-max}{8}
\setval{search-tau-min}{4}
\setval{search-delta}{0,010}
\setval{search-folds}{4}  % folds of training edges for the weight step
\setval{search-alr}{0,001}  % learning rate of the alpha step
\setval{search-weight-decay}{0}
\setval{search-input-drop}{0}
\setval{env-gpu}{RTX 4070 SUPER}  % the run records
\setval{env-gpu-mem}{12}  % GB, the card's specification
\setval{env-python}{3.9.25}
\setval{env-torch}{2.1.2}
\setval{env-cuda}{12.1}
\setval{env-dgl}{2.2.1}
\setval{env-numpy}{1.26.4}
\setval{search-runs}{30}
\setval{search-runs-total}{60}
\setval{search-cands-raw}{55}  % P...P paths of at most 5 hops before merging reversed pairs
\setval{search-epoch-s-4}{6,9}  % median seconds per training epoch
\setval{search-prop-s-4}{37}  % median seconds of feature propagation
\setval{search-run-min-4}{18}  % median minutes per run (epochs + propagation)
\setval{search-fleet-h-4}{9,0}  % GPU hours for the 30 runs
\setval{search-winners-4}{26}  % distinct winners among the 30 runs
\setval{search-delta-alt}{0,005}  % the second tolerance of the vote (sensitivity check)
\setval{search-tie-med-4}{60}  % median tie-group size (delta 0.010)
\setval{search-tie-min-4}{13}
\setval{search-tie-max-4}{87}
\setval{search-tie-seed1-4}{83}  % tie group of seed 1
\setval{search-noise-4}{0,0055}  % score noise of one set
\setval{search-agree-4}{\ensuremath{+0,68}}  % Spearman of two runs on shared sets
\setval{search-xnoise-4}{0,0050}  % gap of a winner re-scored in another run
\setval{search-vote-4}{16}  % distinct vote picks
\setval{search-null-med-4}{21}
\setval{search-null-lo-4}{17}  % 90% range, as reselect_fleet.py
\setval{search-null-hi-4}{26}
\setval{search-null-share-4}{3,0}  % null fleets with as few distinct picks
\setval{search-vote-alt-4}{20}  % distinct vote picks
\setval{search-null-med-alt-4}{26}
\setval{search-null-lo-alt-4}{23}  % 90% range, as reselect_fleet.py
\setval{search-null-hi-alt-4}{29}
\setval{search-null-share-alt-4}{0,0}  % null fleets with as few distinct picks
\setval{search-cons-count-4}{5}  % runs whose vote pick is the consensus
\setval{search-null-fleets}{200}  % shuffled-score fleets in the null
\setval{search-avg-sets-4}{333}  % sets scored in at least 5 runs
\setval{search-se-4}{0,0025}  % standard error of the best averaged set
\setval{search-window-4}{2}  % sets within one SE of the best averaged score
\setval{search-epoch-s-6}{8,1}  % median seconds per training epoch
\setval{search-prop-s-6}{36}  % median seconds of feature propagation
\setval{search-run-min-6}{21}  % median minutes per run (epochs + propagation)
\setval{search-fleet-h-6}{10,7}  % GPU hours for the 30 runs
\setval{search-winners-6}{30}  % distinct winners among the 30 runs
\setval{search-tie-med-6}{78}  % median tie-group size (delta 0.010)
\setval{search-tie-min-6}{50}
\setval{search-tie-max-6}{146}
\setval{search-tie-seed1-6}{72}  % tie group of seed 1
\setval{search-noise-6}{0,0068}  % score noise of one set
\setval{search-agree-6}{\ensuremath{+0,59}}  % Spearman of two runs on shared sets
\setval{search-xnoise-6}{0,0039}  % gap of a winner re-scored in another run
\setval{search-vote-6}{19}  % distinct vote picks
\setval{search-null-med-6}{22}
\setval{search-null-lo-6}{19}  % 90% range, as reselect_fleet.py
\setval{search-null-hi-6}{25}
\setval{search-null-share-6}{11,0}  % null fleets with as few distinct picks
\setval{search-vote-alt-6}{20}  % distinct vote picks
\setval{search-null-med-alt-6}{25}
\setval{search-null-lo-alt-6}{22}  % 90% range, as reselect_fleet.py
\setval{search-null-hi-alt-6}{28}
\setval{search-null-share-alt-6}{1,0}  % null fleets with as few distinct picks
\setval{search-cons-count-6}{5}  % runs whose vote pick is the consensus
\setval{search-avg-sets-6}{169}  % sets scored in at least 5 runs
\setval{search-se-6}{0,0030}  % standard error of the best averaged set
\setval{search-window-6}{32}  % sets within one SE of the best averaged score
\setval{pilot-sets}{210}  % sets scored in each same-seed pilot run, all identical
\setval{search-total-h}{20}  % GPU hours for all 60 runs (epochs + propagation)
\setval{rank-pupup-4}{3,6}  % mean alpha rank of PUPUP (1 = most preferred of 36)
\setval{rank-pupup-6}{3,6}  % mean alpha rank of PUPUP (1 = most preferred of 36)
\setval{rank-puppup-4}{2,8}  % mean alpha rank of PUPPUP (1 = most preferred of 36)
\setval{rank-puppup-6}{4,1}  % mean alpha rank of PUPPUP (1 = most preferred of 36)
\setval{rank-pup-4}{9,1}  % mean alpha rank of PUP (1 = most preferred of 36)
\setval{rank-pup-6}{10,2}  % mean alpha rank of PUP (1 = most preferred of 36)
\setval{rank-pcp-4}{26,6}  % mean alpha rank of PCP (1 = most preferred of 36)
\setval{rank-pcp-6}{28,7}  % mean alpha rank of PCP (1 = most preferred of 36)
\setval{rank-pcccp-4}{22,3}  % mean alpha rank of PCCCP (1 = most preferred of 36)
\setval{rank-pcccp-6}{21,8}  % mean alpha rank of PCCCP (1 = most preferred of 36)
\setval{rank-pcpcp-4}{24,5}  % mean alpha rank of PCPCP (1 = most preferred of 36)
\setval{rank-pcpcp-6}{27,5}  % mean alpha rank of PCPCP (1 = most preferred of 36)
\setval{rank-ppccpp-4}{14,9}  % mean alpha rank of PPCCPP (1 = most preferred of 36)
\setval{rank-ppccpp-6}{6,4}  % mean alpha rank of PPCCPP (1 = most preferred of 36)

% --- Settings: the commands in the seed-1 logs of M5 (HAN, metapath2vec) and of GAT --------------
\setval{han-set-topk}{20}  % --topk
\setval{han-set-hidden}{512}  % --hidden
\setval{han-set-num-heads}{8}  % --num-heads
\setval{han-set-feat-drop}{0,2}  % --feat-drop
\setval{han-set-attn-drop}{0,2}  % --attn-drop
\setval{han-set-lr}{0,001}  % --lr
\setval{han-set-weight-decay}{0}  % --weight-decay
\setval{han-set-batch-size}{100.000}  % --batch-size
\setval{han-set-stage}{150}  % --stage
\setval{han-set-patience}{10}  % --patience
\setval{han-set-min-delta}{0,0005}  % --min-delta
\setval{han-set-beta-warmup}{15}  % --beta-warmup
\setval{mtv-set-dim}{512}  % --dim
\setval{mtv-set-walks-per-node}{5}  % --walks-per-node
\setval{mtv-set-walk-length}{20}  % --walk-length
\setval{mtv-set-window}{3}  % --window
\setval{mtv-set-negatives}{5}  % --negatives
\setval{mtv-set-lr}{0,01}  % --lr
\setval{mtv-set-epochs}{10}  % --epochs
\setval{mtv-set-evals-per-epoch}{4}  % --evals-per-epoch
\setval{mtv-set-patience}{6}  % --patience
\setval{mtv-set-min-delta}{0,0005}  % --min-delta
\setval{mtv-set-batch-walks}{1.024}  % --batch-walks
\setval{gat-set-in-dim}{256}  % --in-dim
\setval{gat-set-hidden}{512}  % --hidden
\setval{gat-set-num-heads}{8}  % --num-heads
\setval{gat-set-layers}{2}  % --layers
\setval{gat-set-feat-drop}{0,2}  % --feat-drop
\setval{gat-set-attn-drop}{0,2}  % --attn-drop
\setval{gat-set-lr}{0,001}  % --lr
\setval{gat-set-weight-decay}{0}  % --weight-decay
\setval{gat-set-batch-size}{100.000}  % --batch-size
\setval{gat-set-stage}{150}  % --stage
\setval{gat-set-patience}{10}  % --patience
\setval{gat-set-min-delta}{0,0005}  % --min-delta
\setval{han-set-head-dim}{64}  % dimensions per attention head
\setval{han-set-sem-dim}{128}  % semantic attention vector (model print in the log)
\setval{han-set-in-dim}{260}  % product input features: LSA title vector and product group
\setval{han-set-semantic-tau}{1}  % softmax temperature

% --- Path sets (HAN repository, arch.py) ---------------------------------------------------------
\setval{runs-total}{68}  % result rows: every model, set and seed
\setval{sets-distinct-paths}{32}  % distinct paths over the 24 sets, all among the candidates
\setval{sets-size-min}{2}  % paths per set, over the 24 sets
\setval{sets-size-max}{6}
\setval{sets-hops-max}{5}  % longest path in the 24 sets, hops
\setval{perm-shuffles}{20.000}  % permutation tests of the Spearman correlations
\setval{set-m1}{PUP PCP}  % amazon_lp_manual_1
\setval{set-m1-n}{2}
\setval{set-m1-hops}{4}
\setval{set-m2}{PP PUP PCP}  % amazon_lp_manual_2
\setval{set-m2-n}{3}
\setval{set-m2-hops}{5}
\setval{set-m3}{PP PUP PPUP}  % amazon_lp_manual_3
\setval{set-m3-n}{3}
\setval{set-m3-hops}{6}
\setval{set-m4}{PCP PCCCP PCPCP}  % amazon_lp_manual_4
\setval{set-m4-n}{3}
\setval{set-m4-hops}{10}
\setval{set-m5}{PP PUP PCP PPUP}  % amazon_lp_manual_5
\setval{set-m5-n}{4}
\setval{set-m5-hops}{8}
\setval{set-m6}{PP PUP PCP PPUP PPP PCCCP}  % amazon_lp_manual_6
\setval{set-m6-n}{6}
\setval{set-m6-hops}{14}
\setval{set-m7}{PP PPP PUP PCP}  % amazon_lp_manual_7
\setval{set-m7-n}{4}
\setval{set-m7-hops}{7}
\setval{set-m8}{PUP PCP PUPUP PCPUP}  % amazon_lp_manual_8
\setval{set-m8-n}{4}
\setval{set-m8-hops}{12}
\setval{set-r4d1}{PCCPPP PCPPUP PPCCPP PPPPP}  % amazon_lp_rand4_d1
\setval{set-r4d1-n}{4}
\setval{set-r4d1-hops}{19}
\setval{set-r4d2}{PCCPCP PCCPP PCPCP PPUP}  % amazon_lp_rand4_d2
\setval{set-r4d2-n}{4}
\setval{set-r4d2-hops}{16}
\setval{set-r4d3}{PCCCP PP PPPPP PUPUP}  % amazon_lp_rand4_d3
\setval{set-r4d3-n}{4}
\setval{set-r4d3-hops}{13}
\setval{set-r4d4}{PCCP PCPPPP PUP PUPPUP}  % amazon_lp_rand4_d4
\setval{set-r4d4-n}{4}
\setval{set-r4d4-hops}{15}
\setval{set-r4d5}{PCCPPP PCPP PCPUPP PPP}  % amazon_lp_rand4_d5
\setval{set-r4d5-n}{4}
\setval{set-r4d5-hops}{15}
\setval{set-r6d1}{PCCCCP PCCCPP PCCPCP PCCPP PCPUP PPUPP}  % amazon_lp_rand6_d1
\setval{set-r6d1-n}{6}
\setval{set-r6d1-hops}{27}
\setval{set-r6d2}{PCCCP PCPPUP PCPUPP PPCPP PPP PPUP}  % amazon_lp_rand6_d2
\setval{set-r6d2-n}{6}
\setval{set-r6d2-hops}{23}
\setval{set-r6d3}{PCCPPP PCPPCP PCPUP PPCPPP PPPPPP PUPUP}  % amazon_lp_rand6_d3
\setval{set-r6d3-n}{6}
\setval{set-r6d3-hops}{28}
\setval{set-r6d4}{PCCPP PPCPPP PPP PPPPP PPPPUP PUPUP}  % amazon_lp_rand6_d4
\setval{set-r6d4-n}{6}
\setval{set-r6d4-hops}{24}
\setval{set-r6d5}{PCCPPP PCPPCP PCPUP PPCPUP PPPPPP PPPUP}  % amazon_lp_rand6_d5
\setval{set-r6d5-n}{6}
\setval{set-r6d5-hops}{28}
\setval{set-f4a1}{PCCPCP PPCCPP PPUPP PUPUP}  % amazon_lp_v3_f4_a1
\setval{set-f4a1-n}{4}
\setval{set-f4a1-hops}{18}
\setval{set-f4a2}{PCCPPP PPCCPP PPPUPP PUPUP}  % amazon_lp_v3_f4_a2
\setval{set-f4a2-n}{4}
\setval{set-f4a2-hops}{19}
\setval{set-f4c}{PCCPP PPUPP PUPPUP PUPUP}  % amazon_lp_v3_f4_c
\setval{set-f4c-n}{4}
\setval{set-f4c-hops}{17}
\setval{set-f6a1}{PCCP PCCPP PPCCPP PPUP PUPPUP PUPUP}  % amazon_lp_v3_f6_a1
\setval{set-f6a1-n}{6}
\setval{set-f6a1-hops}{24}
\setval{set-f6a2}{PCCP PCCPCP PPCCPP PPUP PUPPUP PUPUP}  % amazon_lp_v3_f6_a2
\setval{set-f6a2-n}{6}
\setval{set-f6a2-hops}{25}
\setval{set-f6c}{PCCP PCCPUP PPCCPP PPUPP PUPPUP PUPUP}  % amazon_lp_v3_f6_c
\setval{set-f6c-n}{6}
\setval{set-f6c-hops}{26}

% --- Every set at seed 1: HAN and metapath2vec (results.csv) -------------------------------------
\setval{han-m1-hit-s1}{0,635}
\setval{han-m1-ndcg-s1}{0,398}
\setval{han-m1-fhit-s1}{0,0169}
\setval{han-m1-fndcg-s1}{0,0082}
\setval{han-m1-frec-s1}{0,0305}
\setval{han-m1-val-s1}{0,3072}
\setval{han-m2-hit-s1}{0,650}
\setval{han-m2-ndcg-s1}{0,413}
\setval{han-m2-fhit-s1}{0,0192}
\setval{han-m2-fndcg-s1}{0,0093}
\setval{han-m2-frec-s1}{0,0337}
\setval{han-m2-val-s1}{0,2957}
\setval{han-m3-hit-s1}{0,648}
\setval{han-m3-ndcg-s1}{0,416}
\setval{han-m3-fhit-s1}{0,0203}
\setval{han-m3-fndcg-s1}{0,0102}
\setval{han-m3-frec-s1}{0,0382}
\setval{han-m3-val-s1}{0,2935}
\setval{han-m4-hit-s1}{0,619}
\setval{han-m4-ndcg-s1}{0,375}
\setval{han-m4-fhit-s1}{0,0121}
\setval{han-m4-fndcg-s1}{0,0058}
\setval{han-m4-frec-s1}{0,0244}
\setval{han-m4-val-s1}{0,3194}
\setval{han-m5-hit-s1}{0,665}
\setval{han-m5-ndcg-s1}{0,427}
\setval{han-m5-fhit-s1}{0,0216}
\setval{han-m5-fndcg-s1}{0,0106}
\setval{han-m5-frec-s1}{0,0388}
\setval{han-m5-val-s1}{0,2814}
\setval{han-m6-hit-s1}{0,692}
\setval{han-m6-ndcg-s1}{0,453}
\setval{han-m6-fhit-s1}{0,0257}
\setval{han-m6-fndcg-s1}{0,0128}
\setval{han-m6-frec-s1}{0,0462}
\setval{han-m6-val-s1}{0,2651}
\setval{han-m7-hit-s1}{0,670}
\setval{han-m7-ndcg-s1}{0,433}
\setval{han-m7-fhit-s1}{0,0228}
\setval{han-m7-fndcg-s1}{0,0114}
\setval{han-m7-frec-s1}{0,0423}
\setval{han-m7-val-s1}{0,2834}
\setval{han-m8-hit-s1}{0,663}
\setval{han-m8-ndcg-s1}{0,427}
\setval{han-m8-fhit-s1}{0,0212}
\setval{han-m8-fndcg-s1}{0,0102}
\setval{han-m8-frec-s1}{0,0385}
\setval{han-m8-val-s1}{0,2875}
\setval{han-r4d1-hit-s1}{0,671}
\setval{han-r4d1-ndcg-s1}{0,441}
\setval{han-r4d1-fhit-s1}{0,0267}
\setval{han-r4d1-fndcg-s1}{0,0133}
\setval{han-r4d1-frec-s1}{0,0501}
\setval{han-r4d1-val-s1}{0,2816}
\setval{han-r4d2-hit-s1}{0,662}
\setval{han-r4d2-ndcg-s1}{0,421}
\setval{han-r4d2-fhit-s1}{0,0203}
\setval{han-r4d2-fndcg-s1}{0,0102}
\setval{han-r4d2-frec-s1}{0,0381}
\setval{han-r4d2-val-s1}{0,2865}
\setval{han-r4d3-hit-s1}{0,677}
\setval{han-r4d3-ndcg-s1}{0,443}
\setval{han-r4d3-fhit-s1}{0,0243}
\setval{han-r4d3-fndcg-s1}{0,0124}
\setval{han-r4d3-frec-s1}{0,0459}
\setval{han-r4d3-val-s1}{0,2785}
\setval{han-r4d4-hit-s1}{0,671}
\setval{han-r4d4-ndcg-s1}{0,436}
\setval{han-r4d4-fhit-s1}{0,0227}
\setval{han-r4d4-fndcg-s1}{0,0113}
\setval{han-r4d4-frec-s1}{0,0413}
\setval{han-r4d4-val-s1}{0,2827}
\setval{han-r4d5-hit-s1}{0,660}
\setval{han-r4d5-ndcg-s1}{0,422}
\setval{han-r4d5-fhit-s1}{0,0229}
\setval{han-r4d5-fndcg-s1}{0,0113}
\setval{han-r4d5-frec-s1}{0,0435}
\setval{han-r4d5-val-s1}{0,2922}
\setval{han-r6d1-hit-s1}{0,673}
\setval{han-r6d1-ndcg-s1}{0,436}
\setval{han-r6d1-fhit-s1}{0,0228}
\setval{han-r6d1-fndcg-s1}{0,0113}
\setval{han-r6d1-frec-s1}{0,0403}
\setval{han-r6d1-val-s1}{0,2821}
\setval{han-r6d2-hit-s1}{0,706}
\setval{han-r6d2-ndcg-s1}{0,468}
\setval{han-r6d2-fhit-s1}{0,0280}
\setval{han-r6d2-fndcg-s1}{0,0142}
\setval{han-r6d2-frec-s1}{0,0501}
\setval{han-r6d2-val-s1}{0,2580}
\setval{han-r6d3-hit-s1}{0,694}
\setval{han-r6d3-ndcg-s1}{0,460}
\setval{han-r6d3-fhit-s1}{0,0271}
\setval{han-r6d3-fndcg-s1}{0,0135}
\setval{han-r6d3-frec-s1}{0,0503}
\setval{han-r6d3-val-s1}{0,2708}
\setval{han-r6d4-hit-s1}{0,675}
\setval{han-r6d4-ndcg-s1}{0,443}
\setval{han-r6d4-fhit-s1}{0,0262}
\setval{han-r6d4-fndcg-s1}{0,0133}
\setval{han-r6d4-frec-s1}{0,0485}
\setval{han-r6d4-val-s1}{0,2782}
\setval{han-r6d5-hit-s1}{0,706}
\setval{han-r6d5-ndcg-s1}{0,469}
\setval{han-r6d5-fhit-s1}{0,0282}
\setval{han-r6d5-fndcg-s1}{0,0142}
\setval{han-r6d5-frec-s1}{0,0522}
\setval{han-r6d5-val-s1}{0,2610}
\setval{han-f4a1-hit-s1}{0,635}
\setval{han-f4a1-ndcg-s1}{0,405}
\setval{han-f4a1-fhit-s1}{0,0205}
\setval{han-f4a1-fndcg-s1}{0,0104}
\setval{han-f4a1-frec-s1}{0,0368}
\setval{han-f4a1-val-s1}{0,3056}
\setval{han-f4a2-hit-s1}{0,650}
\setval{han-f4a2-ndcg-s1}{0,421}
\setval{han-f4a2-fhit-s1}{0,0232}
\setval{han-f4a2-fndcg-s1}{0,0116}
\setval{han-f4a2-frec-s1}{0,0407}
\setval{han-f4a2-val-s1}{0,3019}
\setval{han-f4c-hit-s1}{0,658}
\setval{han-f4c-ndcg-s1}{0,430}
\setval{han-f4c-fhit-s1}{0,0233}
\setval{han-f4c-fndcg-s1}{0,0115}
\setval{han-f4c-frec-s1}{0,0441}
\setval{han-f4c-val-s1}{0,2984}
\setval{han-f6a1-hit-s1}{0,666}
\setval{han-f6a1-ndcg-s1}{0,435}
\setval{han-f6a1-fhit-s1}{0,0246}
\setval{han-f6a1-fndcg-s1}{0,0124}
\setval{han-f6a1-frec-s1}{0,0449}
\setval{han-f6a1-val-s1}{0,2885}
\setval{han-f6a2-hit-s1}{0,666}
\setval{han-f6a2-ndcg-s1}{0,435}
\setval{han-f6a2-fhit-s1}{0,0238}
\setval{han-f6a2-fndcg-s1}{0,0121}
\setval{han-f6a2-frec-s1}{0,0448}
\setval{han-f6a2-val-s1}{0,2887}
\setval{han-f6c-hit-s1}{0,666}
\setval{han-f6c-ndcg-s1}{0,435}
\setval{han-f6c-fhit-s1}{0,0249}
\setval{han-f6c-fndcg-s1}{0,0127}
\setval{han-f6c-frec-s1}{0,0451}
\setval{han-f6c-val-s1}{0,2885}
\setval{mtv-m1-hit-s1}{0,600}
\setval{mtv-m1-ndcg-s1}{0,374}
\setval{mtv-m1-fhit-s1}{0,0122}
\setval{mtv-m1-fndcg-s1}{0,0061}
\setval{mtv-m1-frec-s1}{0,0188}
\setval{mtv-m1-val-s1}{0,3401}
\setval{mtv-m2-hit-s1}{0,638}
\setval{mtv-m2-ndcg-s1}{0,424}
\setval{mtv-m2-fhit-s1}{0,0263}
\setval{mtv-m2-fndcg-s1}{0,0129}
\setval{mtv-m2-frec-s1}{0,0459}
\setval{mtv-m2-val-s1}{0,3174}
\setval{mtv-m3-hit-s1}{0,624}
\setval{mtv-m3-ndcg-s1}{0,431}
\setval{mtv-m3-fhit-s1}{0,0446}
\setval{mtv-m3-fndcg-s1}{0,0230}
\setval{mtv-m3-frec-s1}{0,0803}
\setval{mtv-m3-val-s1}{0,3260}
\setval{mtv-m4-hit-s1}{0,095}
\setval{mtv-m4-ndcg-s1}{0,042}
\setval{mtv-m4-fhit-s1}{0,0002}
\setval{mtv-m4-fndcg-s1}{0,0001}
\setval{mtv-m4-frec-s1}{0,0003}
\setval{mtv-m4-val-s1}{0,6931}
\setval{mtv-m5-hit-s1}{0,639}
\setval{mtv-m5-ndcg-s1}{0,427}
\setval{mtv-m5-fhit-s1}{0,0273}
\setval{mtv-m5-fndcg-s1}{0,0136}
\setval{mtv-m5-frec-s1}{0,0452}
\setval{mtv-m5-val-s1}{0,3119}
\setval{mtv-m6-hit-s1}{0,639}
\setval{mtv-m6-ndcg-s1}{0,432}
\setval{mtv-m6-fhit-s1}{0,0315}
\setval{mtv-m6-fndcg-s1}{0,0157}
\setval{mtv-m6-frec-s1}{0,0568}
\setval{mtv-m6-val-s1}{0,3166}
\setval{mtv-m7-hit-s1}{0,633}
\setval{mtv-m7-ndcg-s1}{0,421}
\setval{mtv-m7-fhit-s1}{0,0272}
\setval{mtv-m7-fndcg-s1}{0,0135}
\setval{mtv-m7-frec-s1}{0,0498}
\setval{mtv-m7-val-s1}{0,3188}
\setval{mtv-m8-hit-s1}{0,619}
\setval{mtv-m8-ndcg-s1}{0,392}
\setval{mtv-m8-fhit-s1}{0,0153}
\setval{mtv-m8-fndcg-s1}{0,0075}
\setval{mtv-m8-frec-s1}{0,0248}
\setval{mtv-m8-val-s1}{0,3266}
\setval{mtv-r4d1-hit-s1}{0,553}
\setval{mtv-r4d1-ndcg-s1}{0,360}
\setval{mtv-r4d1-fhit-s1}{0,0179}
\setval{mtv-r4d1-fndcg-s1}{0,0090}
\setval{mtv-r4d1-frec-s1}{0,0353}
\setval{mtv-r4d1-val-s1}{0,4152}
\setval{mtv-r4d2-hit-s1}{0,583}
\setval{mtv-r4d2-ndcg-s1}{0,368}
\setval{mtv-r4d2-fhit-s1}{0,0107}
\setval{mtv-r4d2-fndcg-s1}{0,0050}
\setval{mtv-r4d2-frec-s1}{0,0168}
\setval{mtv-r4d2-val-s1}{0,3733}
\setval{mtv-r4d3-hit-s1}{0,607}
\setval{mtv-r4d3-ndcg-s1}{0,406}
\setval{mtv-r4d3-fhit-s1}{0,0316}
\setval{mtv-r4d3-fndcg-s1}{0,0158}
\setval{mtv-r4d3-frec-s1}{0,0589}
\setval{mtv-r4d3-val-s1}{0,3378}
\setval{mtv-r4d4-hit-s1}{0,631}
\setval{mtv-r4d4-ndcg-s1}{0,419}
\setval{mtv-r4d4-fhit-s1}{0,0263}
\setval{mtv-r4d4-fndcg-s1}{0,0136}
\setval{mtv-r4d4-frec-s1}{0,0436}
\setval{mtv-r4d4-val-s1}{0,3151}
\setval{mtv-r4d5-hit-s1}{0,548}
\setval{mtv-r4d5-ndcg-s1}{0,348}
\setval{mtv-r4d5-fhit-s1}{0,0149}
\setval{mtv-r4d5-fndcg-s1}{0,0074}
\setval{mtv-r4d5-frec-s1}{0,0280}
\setval{mtv-r4d5-val-s1}{0,4075}
\setval{mtv-r6d1-hit-s1}{0,583}
\setval{mtv-r6d1-ndcg-s1}{0,368}
\setval{mtv-r6d1-fhit-s1}{0,0128}
\setval{mtv-r6d1-fndcg-s1}{0,0061}
\setval{mtv-r6d1-frec-s1}{0,0230}
\setval{mtv-r6d1-val-s1}{0,3762}
\setval{mtv-r6d2-hit-s1}{0,622}
\setval{mtv-r6d2-ndcg-s1}{0,409}
\setval{mtv-r6d2-fhit-s1}{0,0226}
\setval{mtv-r6d2-fndcg-s1}{0,0110}
\setval{mtv-r6d2-frec-s1}{0,0393}
\setval{mtv-r6d2-val-s1}{0,3374}
\setval{mtv-r6d3-hit-s1}{0,623}
\setval{mtv-r6d3-ndcg-s1}{0,403}
\setval{mtv-r6d3-fhit-s1}{0,0205}
\setval{mtv-r6d3-fndcg-s1}{0,0100}
\setval{mtv-r6d3-frec-s1}{0,0373}
\setval{mtv-r6d3-val-s1}{0,3237}
\setval{mtv-r6d4-hit-s1}{0,632}
\setval{mtv-r6d4-ndcg-s1}{0,436}
\setval{mtv-r6d4-fhit-s1}{0,0391}
\setval{mtv-r6d4-fndcg-s1}{0,0198}
\setval{mtv-r6d4-frec-s1}{0,0761}
\setval{mtv-r6d4-val-s1}{0,3190}
\setval{mtv-r6d5-hit-s1}{0,607}
\setval{mtv-r6d5-ndcg-s1}{0,394}
\setval{mtv-r6d5-fhit-s1}{0,0187}
\setval{mtv-r6d5-fndcg-s1}{0,0091}
\setval{mtv-r6d5-frec-s1}{0,0340}
\setval{mtv-r6d5-val-s1}{0,3509}
\setval{mtv-f4a1-hit-s1}{0,605}
\setval{mtv-f4a1-ndcg-s1}{0,400}
\setval{mtv-f4a1-fhit-s1}{0,0290}
\setval{mtv-f4a1-fndcg-s1}{0,0150}
\setval{mtv-f4a1-frec-s1}{0,0473}
\setval{mtv-f4a1-val-s1}{0,3442}
\setval{mtv-f4a2-hit-s1}{0,608}
\setval{mtv-f4a2-ndcg-s1}{0,405}
\setval{mtv-f4a2-fhit-s1}{0,0302}
\setval{mtv-f4a2-fndcg-s1}{0,0158}
\setval{mtv-f4a2-frec-s1}{0,0516}
\setval{mtv-f4a2-val-s1}{0,3409}
\setval{mtv-f4c-hit-s1}{0,626}
\setval{mtv-f4c-ndcg-s1}{0,421}
\setval{mtv-f4c-fhit-s1}{0,0364}
\setval{mtv-f4c-fndcg-s1}{0,0193}
\setval{mtv-f4c-frec-s1}{0,0582}
\setval{mtv-f4c-val-s1}{0,3222}
\setval{mtv-f6a1-hit-s1}{0,588}
\setval{mtv-f6a1-ndcg-s1}{0,385}
\setval{mtv-f6a1-fhit-s1}{0,0265}
\setval{mtv-f6a1-fndcg-s1}{0,0142}
\setval{mtv-f6a1-frec-s1}{0,0399}
\setval{mtv-f6a1-val-s1}{0,3528}
\setval{mtv-f6a2-hit-s1}{0,592}
\setval{mtv-f6a2-ndcg-s1}{0,384}
\setval{mtv-f6a2-fhit-s1}{0,0252}
\setval{mtv-f6a2-fndcg-s1}{0,0132}
\setval{mtv-f6a2-frec-s1}{0,0381}
\setval{mtv-f6a2-val-s1}{0,3503}
\setval{mtv-f6c-hit-s1}{0,598}
\setval{mtv-f6c-ndcg-s1}{0,392}
\setval{mtv-f6c-fhit-s1}{0,0266}
\setval{mtv-f6c-fndcg-s1}{0,0141}
\setval{mtv-f6c-frec-s1}{0,0403}
\setval{mtv-f6c-val-s1}{0,3435}

% --- The main comparison: 3 seeds, Welch tests (HAN notes, Part 10) ------------------------------
\setval{han-f4a1-hit-mean}{0,635}
\setval{han-f4a1-hit-sd}{0,003}
\setval{han-m5-hit-mean}{0,664}
\setval{han-m5-hit-sd}{0,002}
\setval{han-4-hit-diff}{\ensuremath{-0,029}}  % f4a1 minus m5
\setval{han-4-hit-rel}{\ensuremath{-4,3}}  % relative difference, %
\setval{han-4-hit-p}{\ensuremath{<}0,001}  % Welch, two-sided
\setval{han-4-hit-pt}{\ensuremath{p<0,001}}
\setval{han-4-hit-diffabs}{0,029}
\setval{han-4-hit-relabs}{4,3}
\setval{han-4-hit-sdx}{10}  % difference in units of the larger seed SD
\setval{han-f4a1-ndcg-mean}{0,405}
\setval{han-f4a1-ndcg-sd}{0,002}
\setval{han-m5-ndcg-mean}{0,426}
\setval{han-m5-ndcg-sd}{0,002}
\setval{han-4-ndcg-diff}{\ensuremath{-0,020}}  % f4a1 minus m5
\setval{han-4-ndcg-rel}{\ensuremath{-4,8}}  % relative difference, %
\setval{han-4-ndcg-p}{\ensuremath{<}0,001}  % Welch, two-sided
\setval{han-4-ndcg-pt}{\ensuremath{p<0,001}}
\setval{han-4-ndcg-diffabs}{0,020}
\setval{han-4-ndcg-relabs}{4,8}
\setval{han-4-ndcg-sdx}{8}  % difference in units of the larger seed SD
\setval{han-f4a1-fhit-mean}{0,0207}
\setval{han-f4a1-fhit-sd}{0,0005}
\setval{han-m5-fhit-mean}{0,0211}
\setval{han-m5-fhit-sd}{0,0004}
\setval{han-4-fhit-diff}{\ensuremath{-0,0004}}  % f4a1 minus m5
\setval{han-4-fhit-rel}{\ensuremath{-2,0}}  % relative difference, %
\setval{han-4-fhit-p}{0,32}  % Welch, two-sided
\setval{han-4-fhit-pt}{\ensuremath{p=0,32}}
\setval{han-4-fhit-diffabs}{0,0004}
\setval{han-4-fhit-relabs}{2,0}
\setval{han-4-fhit-sdx}{1}  % difference in units of the larger seed SD
\setval{han-f4a1-fndcg-mean}{0,0103}
\setval{han-f4a1-fndcg-sd}{0,0002}
\setval{han-m5-fndcg-mean}{0,0104}
\setval{han-m5-fndcg-sd}{0,0002}
\setval{han-4-fndcg-diff}{\ensuremath{-0,0002}}  % f4a1 minus m5
\setval{han-4-fndcg-rel}{\ensuremath{-1,7}}  % relative difference, %
\setval{han-4-fndcg-p}{0,37}  % Welch, two-sided
\setval{han-4-fndcg-pt}{\ensuremath{p=0,37}}
\setval{han-4-fndcg-diffabs}{0,0002}
\setval{han-4-fndcg-relabs}{1,7}
\setval{han-4-fndcg-sdx}{1}  % difference in units of the larger seed SD
\setval{han-f4a1-frec-mean}{0,0376}
\setval{han-f4a1-frec-sd}{0,0007}
\setval{han-m5-frec-mean}{0,0379}
\setval{han-m5-frec-sd}{0,0008}
\setval{han-4-frec-diff}{\ensuremath{-0,0004}}  % f4a1 minus m5
\setval{han-4-frec-rel}{\ensuremath{-1,0}}  % relative difference, %
\setval{han-4-frec-p}{0,57}  % Welch, two-sided
\setval{han-4-frec-pt}{\ensuremath{p=0,57}}
\setval{han-4-frec-diffabs}{0,0004}
\setval{han-4-frec-relabs}{1,0}
\setval{han-4-frec-sdx}{0}  % difference in units of the larger seed SD
\setval{han-f6a1-hit-mean}{0,666}
\setval{han-f6a1-hit-sd}{0,002}
\setval{han-m6-hit-mean}{0,687}
\setval{han-m6-hit-sd}{0,004}
\setval{han-6-hit-diff}{\ensuremath{-0,021}}  % f6a1 minus m6
\setval{han-6-hit-rel}{\ensuremath{-3,0}}  % relative difference, %
\setval{han-6-hit-p}{0,0076}  % Welch, two-sided
\setval{han-6-hit-pt}{\ensuremath{p=0,0076}}
\setval{han-6-hit-diffabs}{0,021}
\setval{han-6-hit-relabs}{3,0}
\setval{han-6-hit-sdx}{5}  % difference in units of the larger seed SD
\setval{han-f6a1-ndcg-mean}{0,435}
\setval{han-f6a1-ndcg-sd}{0,001}
\setval{han-m6-ndcg-mean}{0,449}
\setval{han-m6-ndcg-sd}{0,004}
\setval{han-6-ndcg-diff}{\ensuremath{-0,014}}  % f6a1 minus m6
\setval{han-6-ndcg-rel}{\ensuremath{-3,1}}  % relative difference, %
\setval{han-6-ndcg-p}{0,017}  % Welch, two-sided
\setval{han-6-ndcg-pt}{\ensuremath{p=0,017}}
\setval{han-6-ndcg-diffabs}{0,014}
\setval{han-6-ndcg-relabs}{3,1}
\setval{han-6-ndcg-sdx}{3}  % difference in units of the larger seed SD
\setval{han-f6a1-fhit-mean}{0,0245}
\setval{han-f6a1-fhit-sd}{0,0004}
\setval{han-m6-fhit-mean}{0,0246}
\setval{han-m6-fhit-sd}{0,0011}
\setval{han-6-fhit-diff}{\ensuremath{-0,0001}}  % f6a1 minus m6
\setval{han-6-fhit-rel}{\ensuremath{-0,5}}  % relative difference, %
\setval{han-6-fhit-p}{0,87}  % Welch, two-sided
\setval{han-6-fhit-pt}{\ensuremath{p=0,87}}
\setval{han-6-fhit-diffabs}{0,0001}
\setval{han-6-fhit-relabs}{0,5}
\setval{han-6-fhit-sdx}{0}  % difference in units of the larger seed SD
\setval{han-f6a1-fndcg-mean}{0,0123}
\setval{han-f6a1-fndcg-sd}{0,0001}
\setval{han-m6-fndcg-mean}{0,0122}
\setval{han-m6-fndcg-sd}{0,0007}
\setval{han-6-fndcg-diff}{\ensuremath{+0,0001}}  % f6a1 minus m6
\setval{han-6-fndcg-rel}{\ensuremath{+0,6}}  % relative difference, %
\setval{han-6-fndcg-p}{0,87}  % Welch, two-sided
\setval{han-6-fndcg-pt}{\ensuremath{p=0,87}}
\setval{han-6-fndcg-diffabs}{0,0001}
\setval{han-6-fndcg-relabs}{0,6}
\setval{han-6-fndcg-sdx}{0}  % difference in units of the larger seed SD
\setval{han-f6a1-frec-mean}{0,0444}
\setval{han-f6a1-frec-sd}{0,0014}
\setval{han-m6-frec-mean}{0,0442}
\setval{han-m6-frec-sd}{0,0018}
\setval{han-6-frec-diff}{\ensuremath{+0,0002}}  % f6a1 minus m6
\setval{han-6-frec-rel}{\ensuremath{+0,5}}  % relative difference, %
\setval{han-6-frec-p}{0,88}  % Welch, two-sided
\setval{han-6-frec-pt}{\ensuremath{p=0,88}}
\setval{han-6-frec-diffabs}{0,0002}
\setval{han-6-frec-relabs}{0,5}
\setval{han-6-frec-sdx}{0}  % difference in units of the larger seed SD
\setval{mtv-f4a1-hit-mean}{0,600}
\setval{mtv-f4a1-hit-sd}{0,004}
\setval{mtv-m5-hit-mean}{0,640}
\setval{mtv-m5-hit-sd}{0,002}
\setval{mtv-4-hit-diff}{\ensuremath{-0,040}}  % f4a1 minus m5
\setval{mtv-4-hit-rel}{\ensuremath{-6,3}}  % relative difference, %
\setval{mtv-4-hit-p}{0,0013}  % Welch, two-sided
\setval{mtv-4-hit-pt}{\ensuremath{p=0,0013}}
\setval{mtv-4-hit-diffabs}{0,040}
\setval{mtv-4-hit-relabs}{6,3}
\setval{mtv-4-hit-sdx}{9}  % difference in units of the larger seed SD
\setval{mtv-f4a1-ndcg-mean}{0,395}
\setval{mtv-f4a1-ndcg-sd}{0,005}
\setval{mtv-m5-ndcg-mean}{0,426}
\setval{mtv-m5-ndcg-sd}{0,001}
\setval{mtv-4-ndcg-diff}{\ensuremath{-0,031}}  % f4a1 minus m5
\setval{mtv-4-ndcg-rel}{\ensuremath{-7,4}}  % relative difference, %
\setval{mtv-4-ndcg-p}{0,0071}  % Welch, two-sided
\setval{mtv-4-ndcg-pt}{\ensuremath{p=0,0071}}
\setval{mtv-4-ndcg-diffabs}{0,031}
\setval{mtv-4-ndcg-relabs}{7,4}
\setval{mtv-4-ndcg-sdx}{6}  % difference in units of the larger seed SD
\setval{mtv-f4a1-fhit-mean}{0,0279}
\setval{mtv-f4a1-fhit-sd}{0,0010}
\setval{mtv-m5-fhit-mean}{0,0272}
\setval{mtv-m5-fhit-sd}{0,0002}
\setval{mtv-4-fhit-diff}{\ensuremath{+0,0007}}  % f4a1 minus m5
\setval{mtv-4-fhit-rel}{\ensuremath{+2,5}}  % relative difference, %
\setval{mtv-4-fhit-p}{0,38}  % Welch, two-sided
\setval{mtv-4-fhit-pt}{\ensuremath{p=0,38}}
\setval{mtv-4-fhit-diffabs}{0,0007}
\setval{mtv-4-fhit-relabs}{2,5}
\setval{mtv-4-fhit-sdx}{1}  % difference in units of the larger seed SD
\setval{mtv-f4a1-fndcg-mean}{0,0147}
\setval{mtv-f4a1-fndcg-sd}{0,0004}
\setval{mtv-m5-fndcg-mean}{0,0136}
\setval{mtv-m5-fndcg-sd}{0,0002}
\setval{mtv-4-fndcg-diff}{\ensuremath{+0,0010}}  % f4a1 minus m5
\setval{mtv-4-fndcg-rel}{\ensuremath{+7,6}}  % relative difference, %
\setval{mtv-4-fndcg-p}{0,039}  % Welch, two-sided
\setval{mtv-4-fndcg-pt}{\ensuremath{p=0,039}}
\setval{mtv-4-fndcg-diffabs}{0,0010}
\setval{mtv-4-fndcg-relabs}{7,6}
\setval{mtv-4-fndcg-sdx}{2}  % difference in units of the larger seed SD
\setval{mtv-f4a1-frec-mean}{0,0434}
\setval{mtv-f4a1-frec-sd}{0,0034}
\setval{mtv-m5-frec-mean}{0,0452}
\setval{mtv-m5-frec-sd}{0,0005}
\setval{mtv-4-frec-diff}{\ensuremath{-0,0018}}  % f4a1 minus m5
\setval{mtv-4-frec-rel}{\ensuremath{-3,9}}  % relative difference, %
\setval{mtv-4-frec-p}{0,46}  % Welch, two-sided
\setval{mtv-4-frec-pt}{\ensuremath{p=0,46}}
\setval{mtv-4-frec-diffabs}{0,0018}
\setval{mtv-4-frec-relabs}{3,9}
\setval{mtv-4-frec-sdx}{1}  % difference in units of the larger seed SD
\setval{mtv-f6a1-hit-mean}{0,589}
\setval{mtv-f6a1-hit-sd}{0,002}
\setval{mtv-m6-hit-mean}{0,637}
\setval{mtv-m6-hit-sd}{0,002}
\setval{mtv-6-hit-diff}{\ensuremath{-0,049}}  % f6a1 minus m6
\setval{mtv-6-hit-rel}{\ensuremath{-7,6}}  % relative difference, %
\setval{mtv-6-hit-p}{\ensuremath{<}0,001}  % Welch, two-sided
\setval{mtv-6-hit-pt}{\ensuremath{p<0,001}}
\setval{mtv-6-hit-diffabs}{0,049}
\setval{mtv-6-hit-relabs}{7,6}
\setval{mtv-6-hit-sdx}{20}  % difference in units of the larger seed SD
\setval{mtv-f6a1-ndcg-mean}{0,384}
\setval{mtv-f6a1-ndcg-sd}{0,002}
\setval{mtv-m6-ndcg-mean}{0,431}
\setval{mtv-m6-ndcg-sd}{0,001}
\setval{mtv-6-ndcg-diff}{\ensuremath{-0,047}}  % f6a1 minus m6
\setval{mtv-6-ndcg-rel}{\ensuremath{-10,9}}  % relative difference, %
\setval{mtv-6-ndcg-p}{\ensuremath{<}0,001}  % Welch, two-sided
\setval{mtv-6-ndcg-pt}{\ensuremath{p<0,001}}
\setval{mtv-6-ndcg-diffabs}{0,047}
\setval{mtv-6-ndcg-relabs}{10,9}
\setval{mtv-6-ndcg-sdx}{21}  % difference in units of the larger seed SD
\setval{mtv-f6a1-fhit-mean}{0,0263}
\setval{mtv-f6a1-fhit-sd}{0,0005}
\setval{mtv-m6-fhit-mean}{0,0314}
\setval{mtv-m6-fhit-sd}{0,0006}
\setval{mtv-6-fhit-diff}{\ensuremath{-0,0051}}  % f6a1 minus m6
\setval{mtv-6-fhit-rel}{\ensuremath{-16,3}}  % relative difference, %
\setval{mtv-6-fhit-p}{\ensuremath{<}0,001}  % Welch, two-sided
\setval{mtv-6-fhit-pt}{\ensuremath{p<0,001}}
\setval{mtv-6-fhit-diffabs}{0,0051}
\setval{mtv-6-fhit-relabs}{16,3}
\setval{mtv-6-fhit-sdx}{9}  % difference in units of the larger seed SD
\setval{mtv-f6a1-fndcg-mean}{0,0140}
\setval{mtv-f6a1-fndcg-sd}{0,0004}
\setval{mtv-m6-fndcg-mean}{0,0157}
\setval{mtv-m6-fndcg-sd}{0,0001}
\setval{mtv-6-fndcg-diff}{\ensuremath{-0,0018}}  % f6a1 minus m6
\setval{mtv-6-fndcg-rel}{\ensuremath{-11,1}}  % relative difference, %
\setval{mtv-6-fndcg-p}{0,0084}  % Welch, two-sided
\setval{mtv-6-fndcg-pt}{\ensuremath{p=0,0084}}
\setval{mtv-6-fndcg-diffabs}{0,0018}
\setval{mtv-6-fndcg-relabs}{11,1}
\setval{mtv-6-fndcg-sdx}{5}  % difference in units of the larger seed SD
\setval{mtv-f6a1-frec-mean}{0,0395}
\setval{mtv-f6a1-frec-sd}{0,0011}
\setval{mtv-m6-frec-mean}{0,0565}
\setval{mtv-m6-frec-sd}{0,0003}
\setval{mtv-6-frec-diff}{\ensuremath{-0,0170}}  % f6a1 minus m6
\setval{mtv-6-frec-rel}{\ensuremath{-30,2}}  % relative difference, %
\setval{mtv-6-frec-p}{\ensuremath{<}0,001}  % Welch, two-sided
\setval{mtv-6-frec-pt}{\ensuremath{p<0,001}}
\setval{mtv-6-frec-diffabs}{0,0170}
\setval{mtv-6-frec-relabs}{30,2}
\setval{mtv-6-frec-sdx}{15}  % difference in units of the larger seed SD
\setval{twin-maxdiff}{0,0002}  % largest change of test Hit@10 when the twin edges are left out, any run
\setval{main-hit-sd-max}{0,004}  % largest seed SD of test Hit@10 among the main runs

% --- GAT and popularity --------------------------------------------------------------------------
\setval{gat-hit-mean}{0,787}
\setval{gat-hit-sd}{0,004}
\setval{pop-hit}{0,423}
\setval{gat-ndcg-mean}{0,553}
\setval{gat-ndcg-sd}{0,005}
\setval{pop-ndcg}{0,249}
\setval{gat-fhit-mean}{0,0476}
\setval{gat-fhit-sd}{0,0049}
\setval{pop-fhit}{0,0094}
\setval{gat-fndcg-mean}{0,0246}
\setval{gat-fndcg-sd}{0,0027}
\setval{pop-fndcg}{0,0049}
\setval{gat-frec-mean}{0,0897}
\setval{gat-frec-sd}{0,0082}
\setval{pop-frec}{0,0172}
\setval{gat-val-mean}{0,1780}  % validation BPR, mean of the 3 seeds
\setval{gat-val-sd}{0,0031}

% --- By the user's number of training edges: Hit@10, mean over the seeds run (HAN logs) ----------
\setval{testdeg-d0}{8}  % test edges in the group
\setval{testdeg-d1}{15.831}  % test edges in the group
\setval{testdeg-d2}{27.948}  % test edges in the group
\setval{testdeg-d3}{38.197}  % test edges in the group
\setval{han-m1-hit-d1}{0,571}
\setval{han-m1-hit-d2}{0,629}
\setval{han-m1-hit-d3}{0,666}
\setval{han-m2-hit-d1}{0,577}
\setval{han-m2-hit-d2}{0,646}
\setval{han-m2-hit-d3}{0,683}
\setval{han-m3-hit-d1}{0,576}
\setval{han-m3-hit-d2}{0,640}
\setval{han-m3-hit-d3}{0,683}
\setval{han-m4-hit-d1}{0,572}
\setval{han-m4-hit-d2}{0,619}
\setval{han-m4-hit-d3}{0,639}
\setval{han-m5-hit-d1}{0,596}
\setval{han-m5-hit-d2}{0,662}
\setval{han-m5-hit-d3}{0,693}
\setval{han-m6-hit-d1}{0,622}
\setval{han-m6-hit-d2}{0,686}
\setval{han-m6-hit-d3}{0,715}
\setval{han-m7-hit-d1}{0,605}
\setval{han-m7-hit-d2}{0,668}
\setval{han-m7-hit-d3}{0,698}
\setval{han-m8-hit-d1}{0,597}
\setval{han-m8-hit-d2}{0,658}
\setval{han-m8-hit-d3}{0,694}
\setval{han-r4d1-hit-d1}{0,601}
\setval{han-r4d1-hit-d2}{0,666}
\setval{han-r4d1-hit-d3}{0,705}
\setval{han-r4d2-hit-d1}{0,595}
\setval{han-r4d2-hit-d2}{0,658}
\setval{han-r4d2-hit-d3}{0,693}
\setval{han-r4d3-hit-d1}{0,609}
\setval{han-r4d3-hit-d2}{0,673}
\setval{han-r4d3-hit-d3}{0,709}
\setval{han-r4d4-hit-d1}{0,604}
\setval{han-r4d4-hit-d2}{0,668}
\setval{han-r4d4-hit-d3}{0,701}
\setval{han-r4d5-hit-d1}{0,598}
\setval{han-r4d5-hit-d2}{0,655}
\setval{han-r4d5-hit-d3}{0,689}
\setval{han-r6d1-hit-d1}{0,605}
\setval{han-r6d1-hit-d2}{0,670}
\setval{han-r6d1-hit-d3}{0,703}
\setval{han-r6d2-hit-d1}{0,639}
\setval{han-r6d2-hit-d2}{0,700}
\setval{han-r6d2-hit-d3}{0,739}
\setval{han-r6d3-hit-d1}{0,629}
\setval{han-r6d3-hit-d2}{0,692}
\setval{han-r6d3-hit-d3}{0,724}
\setval{han-r6d4-hit-d1}{0,610}
\setval{han-r6d4-hit-d2}{0,676}
\setval{han-r6d4-hit-d3}{0,702}
\setval{han-r6d5-hit-d1}{0,640}
\setval{han-r6d5-hit-d2}{0,702}
\setval{han-r6d5-hit-d3}{0,736}
\setval{han-f4a1-hit-d1}{0,562}
\setval{han-f4a1-hit-d2}{0,631}
\setval{han-f4a1-hit-d3}{0,668}
\setval{han-f4a2-hit-d1}{0,579}
\setval{han-f4a2-hit-d2}{0,645}
\setval{han-f4a2-hit-d3}{0,683}
\setval{han-f4c-hit-d1}{0,584}
\setval{han-f4c-hit-d2}{0,653}
\setval{han-f4c-hit-d3}{0,694}
\setval{han-f6a1-hit-d1}{0,592}
\setval{han-f6a1-hit-d2}{0,663}
\setval{han-f6a1-hit-d3}{0,700}
\setval{han-f6a2-hit-d1}{0,596}
\setval{han-f6a2-hit-d2}{0,661}
\setval{han-f6a2-hit-d3}{0,699}
\setval{han-f6c-hit-d1}{0,595}
\setval{han-f6c-hit-d2}{0,662}
\setval{han-f6c-hit-d3}{0,698}
\setval{mtv-m1-hit-d1}{0,514}
\setval{mtv-m1-hit-d2}{0,591}
\setval{mtv-m1-hit-d3}{0,641}
\setval{mtv-m2-hit-d1}{0,581}
\setval{mtv-m2-hit-d2}{0,644}
\setval{mtv-m2-hit-d3}{0,657}
\setval{mtv-m3-hit-d1}{0,570}
\setval{mtv-m3-hit-d2}{0,634}
\setval{mtv-m3-hit-d3}{0,639}
\setval{mtv-m4-hit-d1}{0,097}
\setval{mtv-m4-hit-d2}{0,096}
\setval{mtv-m4-hit-d3}{0,093}
\setval{mtv-m5-hit-d1}{0,558}
\setval{mtv-m5-hit-d2}{0,643}
\setval{mtv-m5-hit-d3}{0,671}
\setval{mtv-m6-hit-d1}{0,582}
\setval{mtv-m6-hit-d2}{0,647}
\setval{mtv-m6-hit-d3}{0,654}
\setval{mtv-m7-hit-d1}{0,552}
\setval{mtv-m7-hit-d2}{0,644}
\setval{mtv-m7-hit-d3}{0,659}
\setval{mtv-m8-hit-d1}{0,561}
\setval{mtv-m8-hit-d2}{0,618}
\setval{mtv-m8-hit-d3}{0,643}
\setval{mtv-r4d1-hit-d1}{0,389}
\setval{mtv-r4d1-hit-d2}{0,562}
\setval{mtv-r4d1-hit-d3}{0,615}
\setval{mtv-r4d2-hit-d1}{0,439}
\setval{mtv-r4d2-hit-d2}{0,589}
\setval{mtv-r4d2-hit-d3}{0,638}
\setval{mtv-r4d3-hit-d1}{0,541}
\setval{mtv-r4d3-hit-d2}{0,622}
\setval{mtv-r4d3-hit-d3}{0,624}
\setval{mtv-r4d4-hit-d1}{0,552}
\setval{mtv-r4d4-hit-d2}{0,635}
\setval{mtv-r4d4-hit-d3}{0,662}
\setval{mtv-r4d5-hit-d1}{0,358}
\setval{mtv-r4d5-hit-d2}{0,539}
\setval{mtv-r4d5-hit-d3}{0,633}
\setval{mtv-r6d1-hit-d1}{0,490}
\setval{mtv-r6d1-hit-d2}{0,578}
\setval{mtv-r6d1-hit-d3}{0,625}
\setval{mtv-r6d2-hit-d1}{0,500}
\setval{mtv-r6d2-hit-d2}{0,635}
\setval{mtv-r6d2-hit-d3}{0,663}
\setval{mtv-r6d3-hit-d1}{0,546}
\setval{mtv-r6d3-hit-d2}{0,629}
\setval{mtv-r6d3-hit-d3}{0,650}
\setval{mtv-r6d4-hit-d1}{0,561}
\setval{mtv-r6d4-hit-d2}{0,647}
\setval{mtv-r6d4-hit-d3}{0,650}
\setval{mtv-r6d5-hit-d1}{0,480}
\setval{mtv-r6d5-hit-d2}{0,615}
\setval{mtv-r6d5-hit-d3}{0,654}
\setval{mtv-f4a1-hit-d1}{0,485}
\setval{mtv-f4a1-hit-d2}{0,599}
\setval{mtv-f4a1-hit-d3}{0,648}
\setval{mtv-f4a2-hit-d1}{0,533}
\setval{mtv-f4a2-hit-d2}{0,613}
\setval{mtv-f4a2-hit-d3}{0,635}
\setval{mtv-f4c-hit-d1}{0,535}
\setval{mtv-f4c-hit-d2}{0,629}
\setval{mtv-f4c-hit-d3}{0,662}
\setval{mtv-f6a1-hit-d1}{0,446}
\setval{mtv-f6a1-hit-d2}{0,590}
\setval{mtv-f6a1-hit-d3}{0,648}
\setval{mtv-f6a2-hit-d1}{0,456}
\setval{mtv-f6a2-hit-d2}{0,591}
\setval{mtv-f6a2-hit-d3}{0,649}
\setval{mtv-f6c-hit-d1}{0,466}
\setval{mtv-f6c-hit-d2}{0,599}
\setval{mtv-f6c-hit-d3}{0,653}
\setval{gat-hit-d1}{0,755}
\setval{gat-hit-d2}{0,798}
\setval{gat-hit-d3}{0,793}
\setval{pop-hit-d1}{0,431}
\setval{pop-hit-d2}{0,435}
\setval{pop-hit-d3}{0,412}
\setval{han-4-hit-d1-diff}{\ensuremath{-0,034}}
\setval{han-4-hit-d2-diff}{\ensuremath{-0,031}}
\setval{han-4-hit-d3-diff}{\ensuremath{-0,025}}
\setval{han-6-hit-d1-diff}{\ensuremath{-0,031}}
\setval{han-6-hit-d2-diff}{\ensuremath{-0,023}}
\setval{han-6-hit-d3-diff}{\ensuremath{-0,015}}
\setval{gat-lead-d1}{\ensuremath{+0,117}}  % GAT minus the HAN validation pick
\setval{gat-lead-d2}{\ensuremath{+0,098}}  % GAT minus the HAN validation pick
\setval{gat-lead-d3}{\ensuremath{+0,054}}  % GAT minus the HAN validation pick

% --- All 24 sets at seed 1: families, order, agreement (HAN notes, Part 10) ----------------------
\setval{han-r4-hit-mean}{0,668}
\setval{han-r4-hit-min}{0,660}
\setval{han-r4-hit-max}{0,677}
\setval{han-r4-fhit-mean}{0,0234}
\setval{han-r4-fhit-min}{0,0203}
\setval{han-r4-fhit-max}{0,0267}
\setval{han-4-pos-hit}{thứ 6, 7 và 8}  % places of the searched sets among the 8 candidate-drawn sets
\setval{han-4-pos-val}{thứ 6, 7 và 8}  % places of the searched sets among the 8 candidate-drawn sets
\setval{han-4-pos-fhit}{thứ 3, 4 và 7}  % places of the searched sets among the 8 candidate-drawn sets
\setval{han-r6-hit-mean}{0,691}
\setval{han-r6-hit-min}{0,673}
\setval{han-r6-hit-max}{0,706}
\setval{han-r6-fhit-mean}{0,0264}
\setval{han-r6-fhit-min}{0,0228}
\setval{han-r6-fhit-max}{0,0282}
\setval{han-6-pos-hit}{thứ 6, 7 và 8}  % places of the searched sets among the 8 candidate-drawn sets
\setval{han-6-pos-val}{thứ 6, 7 và 8}  % places of the searched sets among the 8 candidate-drawn sets
\setval{han-6-pos-fhit}{thứ 5, 6 và 7}  % places of the searched sets among the 8 candidate-drawn sets
\setval{mtv-r4-hit-mean}{0,585}
\setval{mtv-r4-hit-min}{0,548}
\setval{mtv-r4-hit-max}{0,631}
\setval{mtv-r4-fhit-mean}{0,0203}
\setval{mtv-r4-fhit-min}{0,0107}
\setval{mtv-r4-fhit-max}{0,0316}
\setval{mtv-4-pos-hit}{thứ 2, 3 và 5}  % places of the searched sets among the 8 candidate-drawn sets
\setval{mtv-4-pos-val}{thứ 2, 4 và 5}  % places of the searched sets among the 8 candidate-drawn sets
\setval{mtv-4-pos-fhit}{thứ 1, 3 và 4}  % places of the searched sets among the 8 candidate-drawn sets
\setval{mtv-r6-hit-mean}{0,613}
\setval{mtv-r6-hit-min}{0,583}
\setval{mtv-r6-hit-max}{0,632}
\setval{mtv-r6-fhit-mean}{0,0227}
\setval{mtv-r6-fhit-min}{0,0128}
\setval{mtv-r6-fhit-max}{0,0391}
\setval{mtv-6-pos-hit}{thứ 5, 6 và 7}  % places of the searched sets among the 8 candidate-drawn sets
\setval{mtv-6-pos-val}{thứ 4, 5 và 7}  % places of the searched sets among the 8 candidate-drawn sets
\setval{mtv-6-pos-fhit}{thứ 2, 3 và 4}  % places of the searched sets among the 8 candidate-drawn sets
\setval{han-r4-vs-manual-diff}{\ensuremath{+0,005}}  % mean of the 5 draws (seed 1) minus the manual set (3 seeds)
\setval{han-r4-vs-manual-rel}{\ensuremath{+0,7}}  % %
\setval{han-r6-vs-manual-diff}{\ensuremath{+0,004}}  % mean of the 5 draws (seed 1) minus the manual set (3 seeds)
\setval{han-r6-vs-manual-rel}{\ensuremath{+0,6}}  % %
\setval{han-valtest-rho}{\ensuremath{-0,94}}
\setval{han-valfull-rho}{\ensuremath{-0,71}}
\setval{han-mtv-rho}{\ensuremath{+0,09}}  % HAN against metapath2vec over the 23 sets (M4 left out)
\setval{han-wins-sampled}{24}  % of 24 sets
\setval{mtv-wins-full}{14}  % of 24 sets
\setval{han-valpick}{r6d2}  % lowest validation BPR among the 24 HAN sets
\setval{mtv-valpick}{m5}  % lowest validation BPR among the 23 metapath2vec sets
\setval{han-4-dense-min}{0,660}
\setval{han-4-dense-max}{0,677}
\setval{han-4-nodense-min}{0,635}
\setval{han-4-nodense-max}{0,658}
\setval{han-6-dense-min}{0,673}
\setval{han-6-dense-max}{0,706}
\setval{han-6-nodense-min}{0,666}
\setval{han-6-nodense-max}{0,666}
\setval{han-4-span}{0,042}  % range of HAN test Hit@10 over all 4-path sets, seed 1
\setval{han-4-nsets}{11}
\setval{han-searched4-mean-s1}{0,648}  % mean of the three searched sets, seed 1
\setval{han-6-span}{0,041}  % range of HAN test Hit@10 over all 6-path sets, seed 1
\setval{han-6-nsets}{9}
\setval{han-searched6-mean-s1}{0,666}  % mean of the three searched sets, seed 1
\setval{searched-with-pupup}{6}  % searched sets holding PUPUP
\setval{searched-with-puppup}{4}  % searched sets holding PUPPUP
\setval{searched-empty-min}{1}  % near-empty category paths per searched set
\setval{searched-empty-max}{3}

% --- What HAN does with the paths: weights at seed 1 against the search (LMSPS part 4b) ----------
\setval{kind-users-paths}{12}
\setval{kind-users-rank}{12,7}
\setval{kind-users-uses}{46}
\setval{kind-users-dropped}{0}
\setval{kind-users-weight}{1,34}
\setval{kind-empty-paths}{6}
\setval{kind-empty-rank}{11,8}
\setval{kind-empty-uses}{25}
\setval{kind-empty-dropped}{16}
\setval{kind-empty-weight}{0,00}
\setval{kind-cat-paths}{10}
\setval{kind-cat-rank}{23,8}
\setval{kind-cat-uses}{22}
\setval{kind-cat-dropped}{3}
\setval{kind-cat-weight}{1,12}
\setval{kind-pp-paths}{4}
\setval{kind-pp-rank}{28,8}
\setval{kind-pp-uses}{16}
\setval{kind-pp-dropped}{2}
\setval{kind-pp-weight}{0,85}
\setval{han-dropped}{21}  % weights below 0.01, 24 sets
\setval{weight-drop}{0,01}  % a path weight below this counts as dropped
\setval{han-empty-kept-min}{0,26}  % largest near-empty-path weight (beta) in a searched set, f4c aside
\setval{han-empty-kept-max}{0,31}
\setval{han-empty-kept-f4c}{0,06}
\setval{weight-pupup}{1,76}  % mean weight of PUPUP as a multiple of an equal share
\setval{uses-pupup}{10}
\setval{weight-puppup}{1,31}  % mean weight of PUPPUP as a multiple of an equal share
\setval{uses-puppup}{5}
\setval{weight-pup}{1,45}  % mean weight of PUP as a multiple of an equal share
\setval{uses-pup}{8}
\setval{weight-pcp}{1,04}  % mean weight of PCP as a multiple of an equal share
\setval{uses-pcp}{7}
\setval{weight-ppccpp}{0,00}  % mean weight of PPCCPP as a multiple of an equal share
\setval{uses-ppccpp}{6}
\setval{pref-cand-rho}{\ensuremath{+0,54}}  % search preference against han, n = 16
\setval{pref-cand-p}{0,035}  % permutation p, 20,000 shuffles, Random(0)
\setval{pref-cand-pt}{\ensuremath{p=0,035}}
\setval{pref-cand-n}{16}
\setval{pref-cand-val-rho}{\ensuremath{-0,61}}  % search preference against val, n = 16
\setval{pref-cand-val-p}{0,017}  % permutation p, 20,000 shuffles, Random(0)
\setval{pref-cand-val-pt}{\ensuremath{p=0,017}}
\setval{pref-cand-val-n}{16}
\setval{pref-cand-full-rho}{\ensuremath{+0,31}}  % search preference against full, n = 16
\setval{pref-cand-full-p}{0,24}  % permutation p, 20,000 shuffles, Random(0)
\setval{pref-cand-full-pt}{\ensuremath{p=0,24}}
\setval{pref-cand-full-n}{16}
\setval{pref-random-rho}{\ensuremath{+0,28}}  % search preference against han, n = 10
\setval{pref-random-p}{0,43}  % permutation p, 20,000 shuffles, Random(0)
\setval{pref-random-pt}{\ensuremath{p=0,43}}
\setval{pref-random-n}{10}
\setval{pref-all-rho}{\ensuremath{+0,18}}  % search preference against han, n = 24
\setval{pref-all-p}{0,39}  % permutation p, 20,000 shuffles, Random(0)
\setval{pref-all-pt}{\ensuremath{p=0,39}}
\setval{pref-all-n}{24}
\setval{pref-mtv-rho}{\ensuremath{+0,29}}  % search preference against m2v, n = 23
\setval{pref-mtv-p}{0,19}  % permutation p, 20,000 shuffles, Random(0)
\setval{pref-mtv-pt}{\ensuremath{p=0,19}}
\setval{pref-mtv-n}{23}

% --- Where the runs differ: Hit@10 by product group, seed 1 (HAN review/evidence/by_group.txt) ---
\setval{bg-pccp-reach-products-pct}{18,9}  % products a P-C-C-P walk reaches
\setval{bg-pccp-reach-test-pct}{27,5}  % test edges whose product it reaches
\setval{bg-4-gap-no}{\ensuremath{-0,023}}  % searched sets minus the sets with a dense category path
\setval{bg-4-gapabs-no}{0,023}
\setval{bg-4-gap-reach}{\ensuremath{-0,010}}  % searched sets minus the sets with a dense category path
\setval{bg-4-gapabs-reach}{0,010}
\setval{bg-4-gap-book}{\ensuremath{-0,026}}  % searched sets minus the sets with a dense category path
\setval{bg-4-gapabs-book}{0,026}
\setval{bg-4-gap-music}{\ensuremath{-0,017}}  % searched sets minus the sets with a dense category path
\setval{bg-4-gapabs-music}{0,017}
\setval{bg-4-gap-dvd}{\ensuremath{-0,009}}  % searched sets minus the sets with a dense category path
\setval{bg-4-gapabs-dvd}{0,009}
\setval{bg-4-gap-video}{\ensuremath{-0,014}}  % searched sets minus the sets with a dense category path
\setval{bg-4-gapabs-video}{0,014}
\setval{bg-4-a1gap-no}{\ensuremath{-0,034}}  % a1 minus the manual set
\setval{bg-4-a1gap-reach}{\ensuremath{-0,019}}  % a1 minus the manual set
\setval{bg-6-gap-no}{\ensuremath{-0,034}}  % searched sets minus the sets with a dense category path
\setval{bg-6-gapabs-no}{0,034}
\setval{bg-6-gap-reach}{\ensuremath{-0,002}}  % searched sets minus the sets with a dense category path
\setval{bg-6-gapabs-reach}{0,002}
\setval{bg-6-gap-book}{\ensuremath{-0,030}}  % searched sets minus the sets with a dense category path
\setval{bg-6-gapabs-book}{0,030}
\setval{bg-6-gap-music}{\ensuremath{-0,034}}  % searched sets minus the sets with a dense category path
\setval{bg-6-gapabs-music}{0,034}
\setval{bg-6-gap-dvd}{\ensuremath{-0,004}}  % searched sets minus the sets with a dense category path
\setval{bg-6-gapabs-dvd}{0,004}
\setval{bg-6-gap-video}{\ensuremath{-0,007}}  % searched sets minus the sets with a dense category path
\setval{bg-6-gapabs-video}{0,007}
\setval{bg-6-a1gap-no}{\ensuremath{-0,031}}  % a1 minus the manual set
\setval{bg-6-a1gap-reach}{\ensuremath{-0,011}}  % a1 minus the manual set
\setval{bg-users}{5.000}  % test users whose full-ranking top 10 is examined
\setval{bg-all-deg}{13,4}  % mean number of training users of a product
\setval{bg-han-deg-med}{74}  % median over the sets, M4 aside
\setval{bg-han-deg-min}{57}
\setval{bg-han-deg-max}{82}
\setval{bg-mtv-deg-med}{6}
\setval{bg-gat-deg}{152}
\setval{bg-pop-deg}{550}
\setval{bg-han-top1-med}{24}  % % of the top 10 from the 1% most-bought
\setval{bg-mtv-top1-max}{0,1}
\setval{bg-gat-top1}{51}
\setval{bg-han-mtv-book}{\ensuremath{+0,015}}  % HAN minus metapath2vec, median over the 23 sets
\setval{bg-han-mtv-music}{\ensuremath{-0,047}}  % HAN minus metapath2vec, median over the 23 sets
\setval{bg-han-mtv-dvd}{\ensuremath{+0,282}}  % HAN minus metapath2vec, median over the 23 sets
\setval{bg-han-mtv-video}{\ensuremath{+0,216}}  % HAN minus metapath2vec, median over the 23 sets

% --- Training and cost (HAN logs; HAN notes, Parts 9-10) -----------------------------------------
\setval{han-runs}{32}
\setval{han-run-min}{9,0}  % median minutes per run
\setval{han-hours}{5,0}  % GPU hours, all runs
\setval{mtv-runs}{32}
\setval{mtv-run-min}{5,2}  % median minutes per run
\setval{mtv-hours}{2,9}  % GPU hours, all runs
\setval{gat-runs}{3}
\setval{gat-run-min}{12,8}  % median minutes per run
\setval{gat-hours}{0,6}  % GPU hours, all runs
\setval{models-hours}{9}  % HAN, metapath2vec and GAT
\setval{han-stop-min}{63}  % epoch (metapath2vec: pass) of the early stop
\setval{han-stop-max}{112}
\setval{gat-stop-min}{130}  % epoch (metapath2vec: pass) of the early stop
\setval{gat-stop-max}{157}
\setval{mtv-stop-min}{1,5}  % epoch (metapath2vec: pass) of the early stop
\setval{mtv-stop-max}{3,0}
\setval{gat-cap}{150}  % epoch cap of GAT
\setval{gat-rerun-seed}{2}  % the GAT seed that hit the cap and was rerun
\setval{gat-rerun-cap}{300}
\setval{gat-rerun-stop}{157}
\setval{han-searched6-min-lo}{11,5}  % minutes per run, the three searched 6-path sets, seed 1
\setval{han-searched6-min-hi}{11,9}
\setval{han-other6-min-lo}{10,1}  % the other 6-path sets (M6, random), seed 1
\setval{han-other6-min-hi}{13,7}
\setval{mtv-searched6-min-lo}{3,4}  % minutes per run, the three searched 6-path sets, seed 1
\setval{mtv-searched6-min-hi}{3,4}
\setval{mtv-other6-min-lo}{3,3}  % the other 6-path sets (M6, random), seed 1
\setval{mtv-other6-min-hi}{10,0}

% --- Repeatability: a same-seed repeat of one HAN run on the same GPU (logs/addition runs) -------
\setval{repeat-han-best}{88}  % best epoch, the same in both runs (HAN, a1 with 4 paths, seed 1)
\setval{repeat-han-val}{0,3056}  % validation BPR, the same in both runs
\setval{repeat-han-hit}{0,6348}  % test Hit@10, the same in both runs
\setval{repeat-han-groupdiff}{0,0001}  % largest Hit@10 difference in a user-degree group
\setval{han-params-min}{27,8}  % millions
\setval{han-params-max}{28,3}  % millions
\setval{han-user-params}{27,4}  % millions: one 512-d vector per user
\setval{han-gain-med}{0,0006}  % validation BPR gained in the last 10 epochs
\setval{han-gain-max}{0,0037}
\setval{gat-gain-min}{0,0007}
\setval{gat-gain-max}{0,0022}
\setval{searched-gain-max}{0,0006}

% --- Side check: HAN with a 5-epoch warm-up, seed 1 (HAN notes, Part 10) -------------------------
\setval{wu5-warmup}{5}  % epochs of path-weight warm-up in the side check (asserted in every log)
\setval{wu5-diff-mean}{\ensuremath{-0,017}}
\setval{wu5-diff-min}{\ensuremath{-0,066}}
\setval{wu5-diff-max}{\ensuremath{+0,007}}
\setval{wu5-dropped}{41}
\setval{wu5-rho}{0,64}
\setval{wu5-4-pos}{thứ 4, 5 và 7}
\setval{wu5-6-pos}{thứ 4, 5 và 7}
\setval{wu5-f6a1-hit}{0,663}
\setval{wu5-m6-hit}{0,644}
\setval{wu5-m6-kept}{3}  % of its 6 paths

% --- From other sources (typed in make_numbers.py, with the source) ------------------------------
\setval{mcauley-raw-reviews}{142,8 triệu}  % McAuley's Amazon product data page, 2014 version (cseweb.ucsd.edu/~jmcauley/datasets/amazon/links.html), read 5 November 2026: raw reviews
\setval{mcauley-dedup-reviews}{83,68 triệu}  % the same page: reviews left after removing duplicate items
\setval{lmsps-t4-gain-min}{1,3}  % LMSPS paper, Table 4: Micro-F1 of HAN* minus HAN, four datasets
\setval{lmsps-t4-gain-max}{2,4}
\setval{dacrema-algos}{18}  % Dacrema et al., RecSys 2019, abstract: neural recommenders analysed
\setval{dacrema-repro}{7}  % the same: reproduced with reasonable effort
\setval{dacrema-beaten}{6}  % the same: often outperformed by simple heuristics
\setval{diffmg-gems-min}{800}  % DiffMG paper, Table 5: GEMS search cost, GPU minutes
\setval{diffmg-gems-max}{2.000}
\setval{diffmg-own-max}{1,6}  % the same table: DiffMG search cost (three runs)

% --- Provenance: the checkouts these numbers were computed from ----------------------------------
\setval{repo-lmsps-commit}{\detokenize{5e32eb6}}  % read from .git/HEAD
\setval{repo-lmsps-branch}{\detokenize{claude/fix-amazon-extraction}}
\setval{repo-han-commit}{\detokenize{2ea02f6}}  % read from .git/HEAD
\setval{repo-han-branch}{\detokenize{claude/han-review}}


% The break keeps "có ... không" on one line; the running head has the title in one piece.
\title[Đánh giá chi phí và hiệu quả của việc tìm kiếm meta-path cho gợi ý sản phẩm dựa trên HAN]{Đánh giá chi phí và hiệu quả của việc tìm kiếm meta-path\texorpdfstring{\\}{ } cho gợi ý sản phẩm dựa trên HAN}
\subtitle{Nghiên cứu có kiểm soát về LMSPS trên Amazon-Meta}
\author{Phan Hoàng Tuấn Anh}
\email{phananh85327@gmail.com}
\affiliation{%
  \department{Khoa Công nghệ Thông tin}
  \institution{Trường Đại học Công Thương Thành phố Hồ Chí Minh (HUIT)}
  \city{Thành phố Hồ Chí Minh}
  \country{Việt Nam}}

\begin{document}

\begin{abstract}
% sections/abstract.tex -- the abstract (paper_v2_vn, translated in L2 from paper_v2's
% sections/abstract.tex, as the sample of the style; terms as in glossary_vn.tex). Every number is
% \val{} from numbers.tex.
HAN và các mô hình tương tự khai thác đồ thị không đồng nhất thông qua các meta-path, vốn thường
được chọn thủ công. Các meta-path do phương pháp tìm kiếm LMSPS tìm được đã giúp HAN phân loại nút
tốt hơn, nhưng việc tìm kiếm như vậy có đáng với chi phí trong bài toán gợi ý hay không thì chưa
được kiểm chứng. Trên Amazon-Meta, chúng tôi so sánh hai cách chọn meta-path cho HAN từ cùng một tập
\val{search-cands} ứng viên: thủ công (cố định từ trước) và bằng LMSPS được điều chỉnh cho bài toán
dự đoán liên kết (tổng hợp từ \val{search-runs} lần chạy), với các tập chọn ngẫu nhiên làm đối
chứng. Mọi tập được thử nghiệm trên HAN và metapath2vec. Các tập tìm được kém các tập thủ công
\val{han-6-hit-relabs}--\val{han-4-hit-relabs}\% theo độ đo đã cố định từ trước và chỉ ngang bằng
khi xếp hạng toàn bộ, sau khoảng \val{search-total-h}~giờ~GPU tìm kiếm. Phần thua kém chủ yếu nằm
ở sách và nhạc, những sản phẩm mà các meta-path danh mục của các tập tìm được không đến được. Đối
chứng ngẫu nhiên cho thấy việc chọn meta-path ít ảnh hưởng đến HAN trên lược đồ này: mọi tập cùng
kích thước chỉ chênh nhau trong phạm vi \val{han-4-span} Hit@10 và các tập chọn ngẫu nhiên đạt kết
quả ít nhất ngang bằng cả hai cách. Với HAN trên lược đồ nhỏ này, kết quả của tìm kiếm tự động khá
gần với thiết kế thủ công, vì vậy lựa chọn giữa hai cách tùy thuộc vào bối cảnh: thiết kế thủ công
tốt hơn khi có chuyên gia làm được việc này trong thời gian ngắn hơn thời gian tìm kiếm, còn tìm
kiếm cho các tập gần tương đương mà không cần kiến thức chuyên môn.

\end{abstract}

% CCS concepts removed in L10 (N19): the paper is not for an ACM venue. To restore them, put back
% \ccsdesc[500]{Information systems~Recommender systems} and
% \ccsdesc[300]{Computing methodologies~Neural networks} here (the heading code above still works).

\keywords{đồ thị không đồng nhất, tìm kiếm meta-path, HAN, hệ thống gợi ý, dự đoán liên kết,
thiết kế meta-path thủ công}

\maketitle
\hypersetup{pdftitle={Đánh giá chi phí và hiệu quả của việc tìm kiếm meta-path cho gợi ý sản phẩm dựa trên HAN}, pdfauthor={Phan Hoàng Tuấn Anh},
  pdfkeywords={đồ thị không đồng nhất, tìm kiếm meta-path, HAN, hệ thống gợi ý, dự đoán liên
  kết, thiết kế meta-path thủ công}}

% sections/intro.tex -- Section 1 (paper_v2_vn, translated in L4 from paper_v2's
% sections/intro.tex; terms as in glossary_vn.tex). The claims stay within paper_v2's; every number
% is \val{} from numbers.tex.
\section{Giới thiệu}
\label{sec:intro}

Các hệ thống gợi ý thường biết nhiều hơn việc ai đã mua gì: sản phẩm thuộc về các danh mục, khách
hàng mua kèm một số sản phẩm với nhau và người dùng nhận xét về chúng. Một đồ thị không đồng nhất
chứa được mọi quan hệ này và meta-path là một cách phổ biến để khai thác nó:
\code{PUP}, sản phẩm--người dùng--sản phẩm, nối hai sản phẩm được cùng một người dùng nhận xét.
HAN \citep{wang2019heterogeneous} biến mỗi meta-path thành một đồ thị và gán trọng số cho các
meta-path bằng cơ chế chú ý. Nhiều mô hình gợi ý được xây dựng trên meta-path
\citep{hu2018leveraging,shi2020hgrec,mu2021fhrec,yan2021attention}. Việc chọn meta-path có ảnh
hưởng: hiệu năng của các mô hình gợi ý dựa trên meta-path
``\foreignlanguage{english}{highly sensitive to the meta-paths they use}'' (rất nhạy với các meta-path mà chúng sử dụng) \citep{ning2022automatic}.

Các meta-path thường được chọn thủ công, một việc đòi hỏi kiến thức chuyên gia và không dùng lại
được từ bộ dữ liệu này sang bộ dữ liệu khác \citep[][Phụ lục~D]{han2020genetic,li2024longrange}.
Hướng phát hiện tự động nhằm đảm nhận việc lựa chọn này (Mục~\ref{sec:related}). LMSPS
\citep{li2024longrange}, một phương pháp tìm kiếm được thực hiện trên một không gian lớn các
meta-path ứng viên, cho biết các meta-path nó tìm được nâng Micro-F1 của HAN thêm
\val{lmsps-t4-gain-min} đến \val{lmsps-t4-gain-max} điểm trên bốn bộ dữ liệu cho bài toán phân loại nút
\citep[Bảng~4]{li2024longrange}. Liệu chúng có giúp ích cho HAN trong bài toán gợi ý hay không
thì, theo hiểu biết của chúng tôi, chưa được kiểm chứng và câu trả lời không hiển nhiên, vì ba lý do.
Thứ nhất, LMSPS đánh giá meta-path bằng mạng riêng mà nó dùng khi tìm kiếm chứ không phải bằng HAN,
nên những meta-path tốt nhất với mạng đó chưa chắc là tốt nhất với HAN. Thứ hai, trong tìm kiếm kiến trúc mạng nơ-ron, các phương pháp tìm kiếm thường không tốt
hơn các lựa chọn ngẫu nhiên \citep{li2019random,yu2020evaluating}. Thứ ba, việc tìm kiếm tốn giờ
GPU, trong khi thiết kế thủ công không cần giờ GPU nào.

Vì vậy, chúng tôi đặt câu hỏi: \emph{so với thiết kế thủ công, tìm kiếm meta-path cho gợi ý
sản phẩm dựa trên HAN có đáng với chi phí không?} Trên Amazon-Meta \citep{leskovec2007viral},
chúng tôi so sánh hai cách chọn meta-path cho HAN từ cùng \val{search-cands} ứng viên.
Cách thứ nhất là thủ công, với các tập được thiết kế với sự hỗ trợ
của AI theo các nguyên tắc nêu rõ và cố định trước khi có bất kỳ kết quả nào. Cách thứ hai là tìm
kiếm bằng LMSPS, được điều chỉnh cho bài toán dự đoán liên kết người dùng--sản phẩm và chạy
\val{search-runs} lần cho mỗi kích thước tập, vì một lần chạy không xếp hạng được các tập tốt nhất
của nó. Để biết việc chọn meta-path ảnh hưởng đến đâu, chúng tôi dùng các tập chọn ngẫu nhiên từ
cùng các ứng viên làm đối chứng. Mỗi tập được huấn luyện trên HAN và metapath2vec
\citep{dong2017metapath2vec} và thử nghiệm trên cùng các cạnh giữ lại. Độ đo chính là Hit@10 lấy
mẫu, cố định từ trước, còn xếp hạng toàn bộ dùng để kiểm tra. Mô hình phổ biến và một GAT không
phân biệt kiểu làm mô hình cơ sở.

Trên lược đồ này, câu trả lời tùy thuộc vào bối cảnh, vì kết quả của tìm kiếm khá gần với thiết kế thủ công.
Theo độ đo đã cố định từ trước, các tập tìm được kém các tập thủ công \val{han-4-hit-relabs}\% và
\val{han-6-hit-relabs}\% trong HAN với bốn và sáu meta-path (\val{han-4-hit-pt} và \val{han-6-hit-pt})
và kém \val{mtv-4-hit-relabs}\% và \val{mtv-6-hit-relabs}\% trong metapath2vec. Trong HAN, chúng ngang bằng
khi xếp hạng toàn bộ, sau khoảng \val{search-total-h}~giờ~GPU tìm kiếm. Vì vậy, thiết kế thủ công là lựa chọn
tốt hơn khi có chuyên gia làm được việc này trong thời gian ngắn hơn thời gian tìm kiếm, còn tìm kiếm cho
các tập gần tương đương khi không có chuyên gia. Các tập tìm được thua chủ yếu ở sách và nhạc, những sản phẩm
mà các meta-path danh mục của chúng không đến được. Đối chứng ngẫu nhiên cho thấy trên lược đồ này việc chọn
meta-path ít ảnh hưởng: trong HAN, các tập chọn ngẫu nhiên đạt kết quả ít nhất ngang bằng cả hai cách. Ngoài
ra, GAT vượt cả hai mô hình dùng meta-path. Có hai giới hạn vẫn cần xem xét. Thứ nhất, khi giai đoạn khởi động
(warm-up) của trọng số meta-path trong HAN ngắn hơn, tập tìm được lại vượt tập thủ công ở sáu meta-path, trong
một phép kiểm tra chỉ dùng một seed (Mục~\ref{sec:res-sensitivity}). Thứ hai, do lược đồ nhỏ, bộ dữ liệu này
không kiểm chứng được trường hợp mà tìm kiếm tự động nhắm tới, đó là những lược đồ quá lớn để thiết kế thủ công
hay chọn ngẫu nhiên.

Các đóng góp của chúng tôi gồm:
\begin{itemize}
  \item Một phiên bản Amazon-Meta đã làm sạch cho bài toán dự đoán liên kết, trong đó việc gộp các
        ấn bản trùng lặp giảm tỷ lệ cạnh thử nghiệm có bản sao trong tập huấn luyện từ
        \val{before-testleak-pct}\% xuống \val{after-testleak-pct}\% (Mục~\ref{sec:data}).
  \item LMSPS được điều chỉnh cho bài toán dự đoán liên kết, kèm các biện pháp chống rò rỉ dữ liệu
        và một cách kết hợp nhiều lần chạy tìm kiếm thành một quyết định
        (Mục~\ref{sec:method}).
  \item Một phép so sánh có kiểm soát giữa các tập meta-path tìm được và thủ công trên HAN và
        metapath2vec, với đối chứng ngẫu nhiên, trong đó các tập và độ đo chính được cố định từ
        trước (Mục~\ref{sec:setup} và~\ref{sec:results}).
  \item Một góc nhìn chi phí--hiệu quả về tìm kiếm so với thiết kế thủ công, với bằng chứng ở mức
        meta-path cho thấy các tập tìm được thua ở đâu (Mục~\ref{sec:results}).
\end{itemize}

% sections/related.tex -- Section 2 (paper_v2_vn, translated in L4 from paper_v2's
% sections/related.tex; terms as in glossary_vn.tex). Every work cited was read for this project;
% each claim was checked against the paper's text.
\section{Các nghiên cứu liên quan}
\label{sec:related}

\paragraph{Gợi ý dựa trên meta-path.} MCRec gán trọng số cho ngữ cảnh dựa trên meta-path
của từng cặp người dùng--sản phẩm bằng cơ chế đồng chú ý (co-attention) \citep{hu2018leveraging}.
Giống HAN \citep{wang2019heterogeneous}, HGRec \citep{shi2020hgrec}, FHRec \citep{mu2021fhrec} và
AMERec \citep{yan2021attention} học một embedding cho mỗi meta-path và hợp nhất
các embedding bằng cơ chế chú ý. Trong khi đó, metapath2vec học embedding của nút từ các bước đi ngẫu
nhiên theo một meta-path \citep{dong2017metapath2vec}. Tất cả các phương pháp này đều coi
meta-path là cho sẵn.

\paragraph{Phát hiện meta-path tự động.} Các phương pháp dựa trên truy vấn tìm $k$ meta-path tốt nhất nối hai nút cho trước \citep{zhu2020topk}. Các phương pháp dựa trên tìm kiếm
học xem một mô hình nên dùng những meta-path nào, hoặc những cấu trúc phong phú hơn. GEMS tiến hóa
các tập cấu trúc meta bằng thuật toán di truyền \citep{han2020genetic} và RMS mở
rộng các tập meta-path bằng học tăng cường \citep{ning2022automatic}. Cả hai đều đánh giá ứng
viên bằng chính mô hình gợi ý. Cách làm này rất chậm vì phải huấn luyện một mô hình gợi ý cho mỗi ứng viên: theo
so sánh của DiffMG, việc tìm kiếm của GEMS mất \val{diffmg-gems-min} đến \val{diffmg-gems-max}~phút~GPU
cho mỗi bộ dữ liệu, trong khi ba lần chạy tìm kiếm của DiffMG chỉ mất nhiều nhất
\val{diffmg-own-max}~phút \citep[Bảng~5]{ding2021diffmg}. Các phương pháp tìm kiếm theo
kiểu DARTS \citep{liu2019darts} giảm chi phí này bằng cách chỉ huấn luyện một lần: mọi ứng viên được
huấn luyện cùng lúc trong một siêu mạng. DiffMG tìm một đồ thị
meta dẫn dắt một mạng nơ-ron đồ thị cho bài toán phân loại nút và dự đoán liên kết
người dùng--sản phẩm, rồi giữ kết quả tốt nhất trong ba lần chạy tìm kiếm theo điểm xác thực
\citep{ding2021diffmg}. LMSPS lấy mẫu dần các meta-path từ một tập ứng viên lớn
\citep{li2024longrange}. Chúng tôi điều chỉnh LMSPS cho bài toán dự đoán liên kết. Khác với GEMS
và RMS, LMSPS đánh giá các tập của nó bằng mạng riêng, một mạng đóng vai trò thay cho HAN
(Mục~\ref{sec:res-weights} và~\ref{sec:conclusion}).

\paragraph{Đánh giá phương pháp tìm kiếm.} LMSPS so sánh các meta-path của nó với các meta-path
gốc của HAN, vốn được chọn thủ công, ngay trong HAN \citep[Bảng~4]{li2024longrange}, nhưng chỉ so
sánh với các meta-path ngẫu nhiên trong mạng riêng của nó, nơi các meta-path ngẫu nhiên đạt điểm thấp
hơn \citep[Bảng~1]{li2024longrange}. Một đối chứng ngẫu nhiên là cần thiết vì trong tìm kiếm kiến
trúc mạng nơ-ron, tìm kiếm ngẫu nhiên kết hợp dừng sớm đạt kết quả ít nhất ngang bằng ENAS, một
phương pháp tìm kiếm hàng đầu, trên hai bộ đánh giá chuẩn \citep{li2019random}. Xét trung bình, các
chiến lược tìm kiếm tiên tiến nhất cho kết quả tương tự chiến lược ngẫu nhiên
\citep{yu2020evaluating}. Trong số các phương pháp tìm kiếm meta-path, RMS vượt một đối chứng ngẫu
nhiên: trong cùng khoảng thời gian, đối chứng này huấn luyện mô hình gợi ý trên các tập chọn ngẫu
nhiên rồi giữ lại tập có điểm xác thực tốt nhất \citep{ning2022automatic}. GEMS và DiffMG không so
sánh với các cấu trúc ngẫu nhiên \citep{han2020genetic,ding2021diffmg}. Chúng tôi so sánh các tập
tìm được với các tập thủ công trên HAN và metapath2vec, nơi các meta-path được sử dụng, với các tập
chọn ngẫu nhiên từ chính các ứng viên của tìm kiếm làm đối chứng.

\paragraph{Đánh giá mô hình gợi ý.} Các mô hình cơ sở đơn giản thường vẫn cho kết quả cạnh tranh:
trong \val{dacrema-algos} mô hình gợi ý nơ-ron gần đây, chỉ kết quả của \val{dacrema-repro} mô hình
là tái hiện được với công sức hợp lý, trong đó \val{dacrema-beaten} mô hình thường có thể bị các
phương pháp heuristic đơn giản vượt qua \citep{dacrema2019progress}. Vì vậy, chúng tôi so sánh thêm
với mô hình phổ biến và một GAT không phân biệt kiểu, vốn là một mô hình cơ sở mạnh trên đồ thị
không đồng nhất \citep{lv2021hgb}. Thứ hạng của các mô hình theo độ đo lấy mẫu có thể khác với thứ
hạng khi xếp hạng toàn bộ \citep{krichene2022sampled}, nên chúng tôi cũng báo cáo kết quả xếp hạng
toàn bộ cho mọi tập meta-path.

% sections/prelim.tex -- Section 3 (paper_v2_vn, translated in L4 from paper_v2's
% sections/prelim.tex; terms as in glossary_vn.tex): definitions as in HAN (its Definitions
% 3.1-3.3) and the task.
\section{Kiến thức cơ sở}
\label{sec:prelim}

\paragraph{Đồ thị không đồng nhất và meta-path.} Theo HAN \citep{wang2019heterogeneous}, một
đồ thị không đồng nhất $G = (V, E)$ đi kèm một ánh xạ kiểu nút $\phi \colon V \to \mathcal{A}$ và
một ánh xạ kiểu cạnh $\psi \colon E \to \mathcal{R}$, trong đó $|\mathcal{A}| + |\mathcal{R}| > 2$.
Đồ thị của chúng tôi có ba kiểu nút, sản phẩm (P, product), danh mục (C, category) và người dùng
(U, user), cùng bốn quan hệ: mua kèm (P--P),
danh mục lá (P--C), cha và con (C--C) và nhận xét tích cực (U--P). Một meta-path
$A_1 \xrightarrow{R_1} A_2 \xrightarrow{R_2} \cdots \xrightarrow{R_l} A_{l+1}$ gồm $l$ bước mô tả
quan hệ hợp thành $R_1 \circ R_2 \circ \cdots \circ R_l$. Chúng tôi viết nó bằng chữ cái đầu của
các kiểu nút. Mọi meta-path ở đây đều bắt đầu và kết thúc tại một sản phẩm: \code{PUP} nối hai sản
phẩm được cùng một người dùng nhận xét, \code{PCP} nối hai sản phẩm trong cùng một danh mục lá.
Láng giềng của một sản phẩm theo một meta-path là những sản phẩm được nối với nó bởi ít nhất một
đường đi theo meta-path đó \citep{wang2019heterogeneous}. Số đường đi như vậy giữa hai sản phẩm có
thể đếm được bằng cách nhân các ma trận kề của các quan hệ trong meta-path.

\paragraph{Bài toán.} Các nhận xét tích cực, tức các cạnh giữa người dùng $\mathcal{U}$ và sản
phẩm $\mathcal{P}$, được chia thành các cạnh huấn luyện, xác thực và thử nghiệm
(Mục~\ref{sec:data}). Đồ thị mà một mô hình thấy chỉ chứa các cạnh huấn luyện. Với một cạnh
giữ lại $(u, p)$, mô hình xếp hạng $p$ trong số các sản phẩm mà $u$ không có cạnh nối tới. Mọi mô
hình cuối cùng đều cho một embedding $e_u$ cho mỗi người dùng và $h_p$ cho mỗi sản phẩm, rồi xếp hạng
theo $e_u^{\top} h_p$ (công thức~\ref{eq:score}). Mục~\ref{sec:setup} trình bày các độ đo.

% sections/data.tex -- Section 4 (paper_v2_vn, translated in L5 from paper_v2's sections/data.tex;
% terms as in glossary_vn.tex). The details that left it in paper_v2 (cluster sizes, the McAuley
% data, the twin check, the title vocabulary, linking every ancestor) are in the supplementary
% document (supplement_vn.tex, N31). The footnote on the McAuley data page, the source of the
% quotation, moved here from the appendix. Every number is \val{} from numbers.tex.
\section{Dữ liệu}
\label{sec:data}

Amazon-Meta\footnote{\url{https://snap.stanford.edu/data/amazon-meta.html}}
\citep{leskovec2007viral} là dữ liệu thu thập từ các trang sản phẩm của Amazon vào mùa hè năm 2006,
gồm \val{raw-records} bản ghi sản phẩm thuộc sách, CD nhạc, DVD và băng video VHS,
\val{raw-reviews} dòng nhận xét và \val{raw-categories} danh mục. Mỗi bản ghi cho biết tiêu đề và
nhóm của sản phẩm, các sản phẩm mà khách hàng mua kèm, các chuỗi danh mục của nó từ gốc đến lá và
các nhận xét về nó, mỗi nhận xét kèm mã khách hàng, ngày, số sao và các lượt bình chọn.

\subsection{Làm sạch dữ liệu thô}

\paragraph{Ấn bản trùng lặp.} Amazon hiển thị các nhận xét về một tác phẩm trên trang của từng ấn
bản và định dạng của nó. Quá trình thu thập đã ghi lại mọi trang. Vì vậy, cùng một dòng nhận
xét xuất hiện dưới nhiều sản phẩm: bản DVD và bản VHS của một bộ phim có chung
\val{example-shared} trong số \val{example-lines} dòng nhận xét của chúng. Tổng cộng,
\val{raw-shared-pct}\% số dòng nhận xét thô cũng xuất hiện dưới một sản phẩm khác. Một bộ lọc chỉ
loại các cặp (khách hàng, sản phẩm) lặp lại vẫn giữ những bản sao này. Trong dữ liệu được xây dựng
theo cách đó (Bảng~\ref{tab:data}, \emph{trước}), \val{before-clustered-pct}\% số sản phẩm nằm
trong \val{before-clusters} cụm ấn bản, còn \val{before-valleak-pct}\% số cạnh xác thực cùng
\val{before-testleak-pct}\% số cạnh thử nghiệm có một bản sao trong tập huấn luyện: cùng một khách
hàng nhận xét một ấn bản khác của cùng tác phẩm. Bản sao làm lộ cạnh giữ lại. Rò rỉ như vậy dẫn
đến việc đánh giá quá cao hiệu năng \citep[\S3.3]{kaufman2012leakage}: một mô hình có thể đạt điểm
cao nhờ tìm ra các bản sao thay vì học sở thích của người dùng. Meta-path tìm bản sao tốt nhất là
meta-path đồng nhận xét ngắn nhất, \code{PUP}, nên sự rò rỉ này ưu ái chính loại meta-path mà một
nghiên cứu về meta-path đem ra so sánh.

Chúng tôi gộp hai sản phẩm khi chúng có chung ít nhất \val{rule-edition-shared} dòng nhận xét giống
hệt nhau (khách hàng, ngày, số sao, lượt bình chọn và lượt bình chọn hữu ích), chiếm ít nhất
\val{rule-edition-overlap-pct}\% số dòng của sản phẩm nhỏ hơn, bỏ qua các dòng xuất hiện dưới hơn
\val{rule-edition-maxrun} sản phẩm. Quy tắc này nối \val{edition-pairs} cặp thành
\val{edition-clusters} cụm. Mỗi cụm giữ thành viên có nhiều dòng nhận xét nhất. Dữ liệu nhận xét
Amazon ra đời sau của \citet{mcauley2015image} cũng có loại trùng lặp này,
``\foreignlanguage{english}{mainly due to near-identical products whose reviews Amazon merges,
e.g.\ VHS and DVD versions of the same movie}'' (chủ yếu do các sản phẩm gần như giống hệt nhau mà
Amazon gộp nhận xét, chẳng hạn các phiên bản VHS và DVD của cùng một bộ phim)\footnote{Trích từ trang
dữ liệu \url{https://cseweb.ucsd.edu/~jmcauley/datasets/amazon/links.html}, truy cập ngày 5 tháng 11
năm 2026.} và đi kèm các tệp đã loại bỏ trùng lặp. Theo hiểu biết của chúng tôi, chưa có bước nào như
vậy được báo cáo cho Amazon-Meta.

\paragraph{Các bước khác.} Chúng tôi loại mọi tài khoản khách hàng có ít nhất
\val{rule-placeholder-repeats} nhận xét cho cùng một sản phẩm (\val{placeholder-ids} tài khoản,
\val{placeholder-lines} dòng nhận xét). Một nhận xét là một tương tác tích cực khi nó có từ
\val{rule-pos-rating} sao trở lên, như \val{raw-pos-pct}\% số dòng thô, trừ khi có ít nhất
\val{rule-helpful-votes} người đọc đã bình chọn cho nó và chưa đến \val{rule-helpful-pct}\% thấy
nó hữu ích, điều này loại \val{raw-unhelpful-pct}\% số nhận xét tích cực. Mỗi người dùng chỉ được
tính một lần cho mỗi tác phẩm. Sau đó, một \val{rule-kcore}-lõi ($k$-core), áp dụng sau khi gộp và
sau quy tắc tài khoản để chỉ đếm các tương tác khác nhau, giữ lại những người dùng và sản phẩm có
ít nhất \val{rule-kcore} tương tác mỗi bên (\val{kcore-iterations} vòng).

% t1_data.tex -- T1: the data before and after the fixes (paper_v2_vn, translated in L3 from
% paper_v2's tables/t1_data.tex; terms as in glossary_vn.tex). Every number is \val{} from
% numbers.tex (make_numbers.py; sources: audit_before_fix.json and data/AMAZON/audit.json).
\begin{table}[htbp]
  \centering
  \small
  \caption{Dữ liệu trước và sau khi gộp các ấn bản trùng lặp và loại các tài khoản có từ
    \val{rule-placeholder-repeats} nhận xét trở lên cho một sản phẩm. Cả hai phiên bản đều giữ các
    nhận xét tích cực và lọc theo \val{rule-kcore}-lõi. Số liệu rà soát trên toàn bộ dữ liệu, seed
    của phép chia~\val{split-seed}.}
  \label{tab:data}
  \begin{tabularx}{\linewidth}{@{}Lrr@{}}
    \toprule
    & Trước & Sau \\
    \midrule
    Sản phẩm / danh mục & \val{before-products} / \val{before-categories}
      & \val{after-products} / \val{after-categories} \\
    Người dùng & \val{before-users} & \val{after-users} \\
    Cạnh người dùng--sản phẩm & \val{before-edges} & \val{after-edges} \\
    Bậc lớn nhất của người dùng & \val{before-maxdeg} & \val{after-maxdeg} \\
    Cụm ấn bản trùng lặp (số sản phẩm trong các cụm) & \val{before-clusters}
      (\val{before-clustered}) & \val{after-clusters} (\val{after-clustered}) \\
    Cạnh lặp lại tương tác của người dùng với một ấn bản khác & \val{before-copyedges-pct}\%
      & \val{after-copyedges-pct}\% \\
    Cạnh xác thực / thử nghiệm có bản sao trong tập huấn luyện & \val{before-valleak-pct}\% /
      \val{before-testleak-pct}\% & \val{after-valleak-pct}\% / \val{after-testleak-pct}\% \\
    Cạnh xác thực / thử nghiệm mà người dùng có cạnh huấn luyện tới một sản phẩm cùng tiêu đề
      (cận trên của số bản trùng còn lại) & \val{before-valtitle-pct}\% /
      \val{before-testtitle-pct}\% & \val{after-valtitle-pct}\% / \val{after-testtitle-pct}\% \\
    \bottomrule
  \end{tabularx}
\end{table}


\paragraph{Tác động.} Việc gộp và quy tắc tài khoản thu nhỏ đồ thị từ \val{before-edges} xuống
\val{after-edges} cạnh người dùng--sản phẩm và giảm tỷ lệ cạnh thử nghiệm có bản sao trong tập huấn
luyện từ \val{before-testleak-pct}\% xuống \val{after-testleak-pct}\% (Bảng~\ref{tab:data}). Một
phép kiểm tra thứ hai, cho các sản phẩm có tập người nhận xét giống hệt nhau, tìm thấy
\val{twin-test} cạnh thử nghiệm có một sản phẩm song sinh như vậy trong tập huấn luyện. Loại bỏ các
cạnh này không làm Hit@10 thử nghiệm của lần chạy nào thay đổi quá \val{twin-maxdiff}.

\subsection{Đồ thị và phép chia}

Đặc trưng của mỗi sản phẩm gồm vectơ
LSA \val{title-dim} chiều của tiêu đề và nhóm sản phẩm. Đặc trưng của mỗi người dùng
chỉ là một vectơ nhị phân \val{user-dim} chiều sinh từ ID khách hàng bằng hàm băm, chứ không dựng từ các nhận xét của người đó, vì chính các nhận xét
này là các cạnh cần dự đoán.

\paragraph{Danh mục lá và meta-path gần rỗng.} Sản phẩm chỉ liên kết với các danh mục lá của nó.
Nếu liên kết thêm mọi danh mục cấp trên, \val{cat-all-largest-pct}\% số sản phẩm sẽ rơi vào cùng
một danh mục. Tuy nhiên, một bước đơn lẻ trên cây phân cấp hiếm khi dẫn
tới một sản phẩm khác: \code{PCCP} nối một sản phẩm với một sản phẩm khác cho
\val{pccp-reach-dvd-pct}\% số DVD và \val{pccp-reach-video-pct}\% số video, nhưng chỉ cho
\val{pccp-reach-book-pct}\% số sách và \val{pccp-reach-music-pct}\% số sản phẩm nhạc. Chúng tôi gọi
sáu meta-path ứng viên có một bước như vậy là \emph{gần rỗng}: \code{PCCP}, \code{PCCPP},
\code{PCCPPP}, \code{PCCPCP}, \code{PCCPUP} và \code{PPCCPP}. Trong các đồ thị của HAN, mỗi meta-path này để \val{empty-isolated-min-pct}--\val{empty-isolated-max-pct}\% số sản phẩm không có láng
giềng nào, trong khi \code{PCCCP}, vốn đến được các danh mục anh em qua danh mục cha của chúng, để
lại \val{pcccp-graph-isolated} sản phẩm. Sách và nhạc chiếm \val{group-bookmusic-pct}\% số sản
phẩm và \val{test-bookmusic-pct}\% số cạnh thử nghiệm, nên các meta-path này sẽ trở lại ở
Mục~\ref{sec:results}.

\paragraph{Phép chia.} Chúng tôi chia \val{after-edges} cạnh người dùng--sản phẩm một lần, ngẫu
nhiên, thành \val{split-train} cạnh huấn luyện, \val{split-val} cạnh xác thực và \val{split-test}
cạnh thử nghiệm (80/10/10, seed của phép chia~\val{split-seed}). Mọi đồ thị và mọi bước đi ngẫu
nhiên chỉ dùng các cạnh huấn luyện. Các cạnh xác thực dẫn dắt việc tìm kiếm và mọi lựa chọn trong
quá trình huấn luyện, còn các cạnh thử nghiệm chỉ dùng để báo cáo kết quả. Phép chia bỏ qua ngày
nhận xét, một hạn chế mà chúng tôi sẽ quay lại ở Mục~\ref{sec:conclusion}.

% sections/method.tex -- Section 5 (paper_v2_vn, translated in L5 from paper_v2's
% sections/method.tex; terms as in glossary_vn.tex). Every number is \val{} from numbers.tex.
% Figure: figures/f2_ties.tex (F2). F1 (the overview), the changes from the base papers with their
% reasons and the search's details are in the supplementary document (supplement_vn.tex, N31). The
% key settings are in Section 6.
\section{Phương pháp}
\label{sec:method}

\subsection{Tìm kiếm meta-path cho dự~đoán~liên~kết}
\label{sec:search}

\paragraph{LMSPS\@.} Trước hết, LMSPS \citep{li2024longrange} cho mỗi meta-path ứng
viên một ma trận đặc trưng trên các nút đích: các đặc trưng gốc ở đầu kia của meta-path được lấy
trung bình qua từng bước dọc theo meta-path. Một siêu mạng gán cho mỗi ứng viên $k$ một trọng số
$\alpha_k$. Ở mỗi epoch siêu mạng này hợp nhất $M$ meta-path lấy mẫu từ một nhóm ứng viên với trọng số
Gumbel-softmax tính từ $\alpha$. Nhóm ứng viên thu hẹp dần còn $V$ meta-path có $\alpha$ lớn
nhất. Trọng số của mạng được học trên dữ liệu huấn luyện và $\alpha$ trên tập dữ liệu xác thực. Cuối cùng, các tập $M$ meta-path từ
nhóm ứng viên cuối được tính điểm trên tập dữ liệu xác thực. Tập tốt nhất là kết quả. Chúng tôi bắt
đầu từ mã nguồn của các tác giả\footnote{\url{https://github.com/JHL-HUST/LMSPS}}, vốn được xây
dựng cho bài toán phân loại nút. Chúng tôi thay đổi mã nguồn cho bài toán dự đoán liên kết.

\paragraph{Điểm và hàm mất mát.} Embedding $h_p$ của một sản phẩm là tổng có trọng số của các đặc
trưng meta-path của nó, mỗi đặc trưng qua một ánh xạ tuyến tính riêng, tiếp theo là PReLU và
dropout. Mỗi người dùng có một embedding tự do $e_u$ \val{search-hidden} chiều. Một cặp được tính điểm
bằng
\begin{equation}
  s(u, p) = e_u^{\top} h_p ,
  \label{eq:score}
\end{equation}
và được huấn luyện với hàm mất mát BPR \citep{rendle2009bpr}, $-\log \sigma\bigl(s(u, p^{+}) -
s(u, p^{-})\bigr)$, trên các cạnh huấn luyện $(u, p^{+})$ và các sản phẩm $p^{-}$ lấy mẫu ngẫu nhiên
từ những sản phẩm mà $u$ không có cạnh huấn luyện nối tới. HAN sẽ dùng cùng thang điểm và hàm mất mát này.

\paragraph{Ứng viên.} Các ứng viên là những meta-path dài từ 1 đến \val{search-hops} bước bắt đầu và
kết thúc tại một sản phẩm, vì HAN cần một đồ thị sản phẩm cho mỗi meta-path:
\val{search-cands-raw} meta-path, hay \val{search-cands} khi mỗi meta-path được gộp với meta-path
ngược của nó cho cùng một đồ thị sau khi đối xứng hóa.

\paragraph{Biện pháp chống rò rỉ.} Các đặc trưng được tính trên đồ thị huấn luyện, vốn chứa cả chính
các cạnh huấn luyện. Vì vậy, với một cạnh huấn luyện, đặc trưng meta-path qua người dùng của sản phẩm
đã chứa sẵn người dùng của cạnh đó, điều mà một cạnh giữ lại không bao giờ có. Để tránh điều này,
chúng tôi học $\alpha$ trên các cạnh xác thực. Ở bước cập nhật trọng số, chúng tôi tính đặc trưng
meta-path qua người dùng \val{search-folds} lần, mỗi lần bỏ đi một phần tư khác nhau của các cạnh
huấn luyện, rồi đọc đặc trưng của mỗi cạnh từ phiên bản không chứa cạnh đó.

\paragraph{Lịch huấn luyện và cách chọn.} Một tập gồm $M = 4$ hoặc $M = 6$ meta-path. Nhóm ứng
viên cuối gồm $V = \val{search-pool}$ meta-path, để có thể tính điểm cả \val{search-sets} tập của nhóm. Một lần chạy
kéo dài \val{search-epochs} epoch, \val{search-warmup} epoch đầu với toàn bộ nhóm ứng viên. LMSPS
tính điểm các tập lấy mẫu với nhiễu Gumbel mới, mà trên dữ liệu của chúng tôi làm điểm của một tập
thay đổi nhiều hơn chênh lệch giữa các tập. Chúng tôi tính điểm mỗi tập một lần, không có nhiễu,
bằng BPR xác thực của nó so với \val{split-negatives} mẫu âm cố định của mỗi cạnh xác thực. Mỗi lần
chạy cũng ghi lại \emph{nhóm đồng hạng} của nó, tức các tập nằm trong phạm vi
$\delta = \val{search-delta}$ so với tập có điểm tốt nhất.

\subsection{Một quyết định từ nhiều lần chạy}
\label{sec:fleet}

Một lần chạy không xếp hạng được các tập tốt nhất của nó (Hình~\ref{fig:ties}): nhóm đồng hạng
của nó chứa trung vị \val{search-tie-med-4} trong \val{search-sets} tập với $M = 4$ meta-path và
\val{search-tie-med-6} với $M = 6$ meta-path. Qua \val{search-runs} lần chạy cho mỗi kích thước, có
\val{search-winners-4} và \val{search-winners-6} tập thắng khác nhau. Vì vậy, chúng tôi chạy tìm
kiếm \val{search-runs} lần cho mỗi kích thước, với seed từ 1 đến \val{search-runs} trên cùng dữ liệu
và phép chia, rồi đưa ra một quyết định từ tất cả các lần chạy.

% f2_ties.tex -- F2: one run cannot rank its top sets (paper_v2_vn, translated in L3 from
% paper_v2's figures/f2_ties.tex). Data: f2_scores.dat and f2_picks.dat, written by
% make_numbers.py from the Fleet v3 records and re-selections.
\begin{figure}[htbp]
  \centering
  \begin{tikzpicture}
    \begin{axis}[
        name=scores, width=\linewidth, height=3.9cm,
        xlabel={Tập, tốt nhất trước}, ylabel={BPR xác thực\\vượt tập tốt nhất},
        ylabel style={align=center},
        xmin=0, xmax=211, ymin=0, ymax=0.125, xtick={1,50,100,150,210},
        yticklabel style={/pgf/number format/fixed, /pgf/number format/precision=2},
        scaled y ticks=false, tick label style={font=\scriptsize}, label style={font=\footnotesize},
        title={(a) Một lần chạy (seed 1): \val{search-sets} tập của nó},
        title style={font=\footnotesize},
        legend style={font=\scriptsize, at={(0.03,0.97)}, anchor=north west, draw=none,
          fill=white, fill opacity=0.85, text opacity=1},
        legend cell align=left, axis line style={black!60}, tick style={black!60}]
      \fill[cSearch!10] (axis cs:0,0) rectangle (axis cs:211,0.010);
      \addplot[cSearch, line width=1pt] table[x=rank, y=gapfour] {figures/f2_scores.dat};
      \addlegendentry{4 meta-path: \val{search-tie-seed1-4} tập trong phạm vi $\delta$}
      \addplot[cSearch!55, line width=1pt, densely dashed] table[x=rank, y=gapsix]
        {figures/f2_scores.dat};
      \addlegendentry{6 meta-path: \val{search-tie-seed1-6} tập trong phạm vi $\delta$}
      \draw[black!55, dashed] (axis cs:0,0.010) -- (axis cs:211,0.010)
        node[pos=0.03, above right, font=\scriptsize, text=black!75]
        {dung sai $\delta = \val{search-delta}$};
    \end{axis}
    \begin{axis}[
        at={(scores.below south west)}, anchor=above north west, yshift=-0.35cm,
        width=\linewidth, height=3.3cm,
        ybar, bar width=9pt, ymin=0, ymax=31, ylabel={Số tập khác nhau},
        symbolic x coords={4,6}, xtick=data, xticklabels={4 meta-path, 6 meta-path},
        enlarge x limits=0.5,
        tick label style={font=\scriptsize}, label style={font=\footnotesize},
        title={(b) Lựa chọn qua \val{search-runs} lần chạy}, title style={font=\footnotesize},
        legend style={font=\scriptsize, at={(0.5,-0.24)}, anchor=north, draw=none,
          legend columns=1}, legend cell align=left,
        axis line style={black!60}, tick style={black!60}]
      \addplot[fill=cSearch, draw=white] table[x=size, y=own] {figures/f2_picks.dat};
      \addlegendentry{tập thắng riêng của từng lần chạy}
      \addplot[fill=cSearch!45, draw=white] table[x=size, y=vote] {figures/f2_picks.dat};
      \addlegendentry{bỏ phiếu qua các lần chạy}
      \addplot[fill=cOther!45, draw=white, error bars/.cd, y dir=both, y explicit,
          error bar style={black!70, line width=0.6pt}]
        table[x=size, y=nullmed, y error plus=errhi, y error minus=errlo] {figures/f2_picks.dat};
      \addlegendentry{bỏ phiếu trên điểm xáo trộn (trung vị, khoảng 90\%)}
    \end{axis}
  \end{tikzpicture}
  \Description{Hai khung, khung này nằm trên khung kia. Trên: điểm xác thực của 210 tập meta-path của một lần chạy tìm kiếm, sắp từ tốt nhất đến kém nhất, với nhiều tập nằm trong dung sai so với tập tốt nhất. Dưới: các cột cho bốn và sáu meta-path, đếm số tập khác nhau mà 30 lần chạy chọn, cho tập thắng riêng của từng lần chạy, cho bỏ phiếu qua các lần chạy và cho bỏ phiếu trên điểm xáo trộn.}
  \caption{Một lần chạy tìm kiếm không xếp hạng được các tập đứng đầu của nó. (a) Điểm xác thực
    của \val{search-sets} tập mà một lần chạy tính điểm, đã sắp xếp: nhiều tập nằm trong dung sai
    $\delta$ so với tập tốt nhất. (b) Qua \val{search-runs} lần chạy, các tập thắng hầu như đều khác
    nhau. Bỏ phiếu cho mức đồng thuận cao hơn, nhưng chỉ cao hơn đôi chút so với khi bỏ phiếu trên
    điểm xáo trộn.}
  \label{fig:ties}
\end{figure}


\paragraph{Điểm trung bình.} Các lần chạy tính điểm những tập khác nhau, vì nhóm ứng viên cuối của
chúng khác nhau, nhưng xếp hạng giống nhau những tập mà chúng cùng có (hệ số Spearman trung vị
\val{search-agree-4} với $M = 4$ và \val{search-agree-6} với $M = 6$).
Chúng tôi khớp giá trị mất mát xác thực của tập $s$ trong lần chạy $r$ dưới dạng $a_r + b_s$ bằng
bình phương tối thiểu và so sánh \val{search-avg-sets-4} và \val{search-avg-sets-6} tập được tính
điểm trong ít nhất năm lần chạy. Sau đó, chúng tôi áp dụng quy tắc một sai số chuẩn: các tập có $b_s$
cách giá trị của tập tốt nhất không quá một sai số chuẩn (\val{search-se-4} với $M = 4$ meta-path và
\val{search-se-6} với $M = 6$ meta-path) tạo thành một cửa sổ, trong đó tập đơn giản nhất theo tổng
số bước được ưu tiên, vì các meta-path ngắn hơn cũng tạo ra đồ thị HAN thưa hơn. Với $M = 4$
meta-path, cửa sổ này chứa \val{search-window-4} tập: \code{a1} (tập tốt nhất) và \code{a2}. Cả hai
đều được chuyển cho các mô hình. Với $M = 6$ meta-path, nó chứa \val{search-window-6} tập gần ngang
nhau, một vùng phẳng mà tìm kiếm không xếp hạng được. Chúng tôi lấy hai tập đơn giản nhất, bằng
số tập với bốn meta-path. Chúng tôi cũng đã ghi lại lý do cho quyết định này.

\paragraph{Bỏ phiếu.} Một quy tắc thứ hai cho mỗi lần chạy chọn, từ nhóm đồng hạng của nó, tập mà
các meta-path của nó xuất hiện thường xuyên nhất trong các tập thắng của những lần chạy khác. Tập
đồng thuận \code{c} gồm $M$ meta-path phổ biến nhất trong các lựa chọn. Vì bản thân cách bỏ phiếu
đã tạo ra một mức đồng thuận nhất định, chúng tôi so sánh nó với \val{search-null-fleets} loạt chạy
có điểm bị xáo trộn trong từng lần chạy (Hình~\ref{fig:ties}b). Mức đồng
thuận giữa các lần chạy cao hơn mức ngẫu nhiên, nhưng không nhiều: \val{search-null-share-4}\% và
\val{search-null-share-6}\% số loạt chạy xáo trộn có số lựa chọn khác nhau ít như vậy.

\paragraph{Những gì được chuyển cho các mô hình.} Mỗi kích thước chuyển giao \code{a1}, \code{a2}
và \code{c} (Bảng~\ref{tab:sets}). Tập \code{a1} là đối tượng của phép so sánh chính với các tập thủ
công, đã được cố định từ trước. Ngưỡng dung sai $\delta$ được cố định trước các lần chạy cuối, dựa
trên nhiễu giữa các lần chạy trước đó. Các điểm trung bình không phụ thuộc vào nó. Một $\delta$
nhỏ hơn cũng không làm thay đổi \code{a1} lẫn \code{a2}.

\subsection{Các mô hình dùng meta-path}
\label{sec:models}

\paragraph{HAN\@.} HAN \citep{wang2019heterogeneous} biến mỗi meta-path thành một đồ thị trên
các nút đích, chạy cơ chế chú ý mức nút, tức một tầng chú ý đồ thị trên từng đồ thị rồi hợp nhất
các embedding theo từng meta-path bằng cơ chế chú ý mức ngữ nghĩa, tức một trọng số $\beta$ cho
mỗi meta-path dùng chung cho mọi nút. Chúng tôi giữ cả hai mức và thay đổi ba điểm.
\emph{Đồ thị:} đồ thị sản phẩm của mỗi meta-path giữ \val{han-set-topk} láng giềng mạnh nhất của
mỗi sản phẩm theo số đường đi (walk count) sau mỗi bước, vì nếu không cắt tỉa, các meta-path có hai
bước qua người dùng sẽ nối một sản phẩm với phần lớn kho sản phẩm. Đồ thị là nhị phân, đối xứng và
cho mỗi sản phẩm một cạnh tự nối. \emph{Người dùng:} HAN chỉ mã hóa các sản phẩm. Mỗi
người dùng có một embedding tự do và một cặp được tính điểm bằng công thức~\eqref{eq:score} với BPR.
Các embedding người dùng chiếm \val{han-user-params}~triệu trong số
\val{han-params-min}--\val{han-params-max}~triệu tham số của HAN. \emph{Trọng số meta-path:} tầng
chú ý của mọi meta-path bắt đầu từ cùng trọng số, còn $\beta$ được giữ bằng nhau trong
\val{han-set-beta-warmup} epoch đầu. Nếu các tầng chú ý được khởi tạo ngẫu nhiên khác nhau cho từng meta-path, trọng số mà HAN gán cho
một meta-path sẽ phụ thuộc vào vị trí của nó trong danh sách meta-path. Không có giai đoạn khởi động, một meta-path đã chiếm toàn bộ
trọng số chỉ sau vài epoch.

\paragraph{metapath2vec.} metapath2vec \citep{dong2017metapath2vec} học embedding của nút, không
dùng đặc trưng, bằng một mô hình skip-gram trên các bước đi ngẫu nhiên theo một meta-path.
Chúng tôi đi mọi meta-path của một tập cùng nhau, mỗi meta-path như một khuôn mẫu lặp lại
(\code{PUP} cho P--U--P--U\ldots), bắt đầu từ mọi sản phẩm, rồi tính điểm các embedding đã cố định
bằng công thức~\eqref{eq:score}. Người dùng chỉ được tiếp cận qua các meta-path có \code{U}, nên
một tập không có meta-path nào như vậy sẽ để người dùng không được huấn luyện.

\paragraph{GAT và mô hình phổ biến.} Làm đối chứng không dùng meta-path, một mạng chú ý đồ thị
\citep{velickovic2018graph} được huấn luyện với BPR, chạy trên hợp của mọi quan hệ, bỏ qua kiểu nút
ngoại trừ một phép chiếu đầu vào cho mỗi kiểu: đây là mô hình cơ sở đồng nhất của bộ đánh giá chuẩn
HGB \citep{lv2021hgb}. Người dùng và sản phẩm cùng từ một mạng, nên một người dùng được dựng từ
các sản phẩm mà người đó đã nhận xét. Mỗi bước huấn luyện bỏ các cạnh dương của chính nó khỏi đồ
thị mà nó chạy trên, vì một cạnh giữ lại không bao giờ nằm trong đồ thị. Mô hình phổ biến tính
điểm một sản phẩm bằng số người dùng huấn luyện của nó (bằng điểm thì xếp ngẫu nhiên) và đưa ra
cùng một bảng xếp hạng cho mọi người dùng.

% sections/setup.tex -- Section 6 (paper_v2_vn, translated in L6 from paper_v2's sections/setup.tex;
% terms as in glossary_vn.tex). The Thiết lập paragraph gives the key settings (the supplementary
% document, supplement_vn.tex, has them all since N31), Table~\ref{tab:cost} the search's run
% times. Every number is \val{} from numbers.tex. Table: tables/t3_sets.tex (T3).
\section{Thiết lập thực nghiệm}
\label{sec:setup}

\paragraph{Các tập meta-path.} Chúng tôi so sánh các tập tìm được với các tập thủ công, cùng các
tập chọn ngẫu nhiên làm đối chứng, tổng cộng 24 tập (Bảng~\ref{tab:sets}).
\begin{itemize}
  \item \emph{Tìm được:} với mỗi kích thước tập, \code{a1}, \code{a2} và \code{c}
        (Mục~\ref{sec:fleet}).
  \item \emph{Thủ công:} M1--M8, thiết kế theo bốn nguyên tắc nêu rõ.
        (1) Mỗi loại bằng chứng một meta-path (mua kèm \code{PP}, đồng nhận xét \code{PUP}, chung
        danh mục \code{PCP}). (2) Ưu tiên meta-path ngắn. (3) Tránh meta-path gần rỗng (Mục~\ref{sec:data}), nên
        \code{PCCCP}, vốn đến được các danh mục anh em, được dùng thay cho \code{PCCP}. (4) Có các tập chính cùng kích thước
        với các tập tìm được, M5 và M6. Sáu tập còn lại, mỗi tập thử một ý tưởng, chẳng hạn chỉ hành vi (M3) hoặc chỉ
        danh mục (M4).
  \item \emph{Ngẫu nhiên (đối chứng):} năm tập cho mỗi kích thước, chọn đều và không lặp lại từ
        chính \val{search-cands} ứng viên của tìm kiếm, để biết việc chọn meta-path ảnh hưởng đến
        đâu.
\end{itemize}
Các tập thủ công và các tập chọn ngẫu nhiên được cố định trước khi có bất kỳ kết quả tìm kiếm hay
mô hình nào trên dữ liệu đã làm sạch. Các tập tìm được thì cố định trước mọi lần chạy HAN. Cả
\val{sets-distinct-paths} meta-path khác nhau trong 24 tập đều nằm trong \val{search-cands} ứng
viên, nên phương pháp tìm kiếm đã có thể chọn bất kỳ meta-path nào trong số đó.

% t3_sets.tex -- T3: the 24 path sets (paper_v2_vn, translated in L3 from paper_v2's
% tables/t3_sets.tex; HAN repository, arch.py, via numbers.tex). Since N31 it also lists the ten
% random draws, which were a table of their own in the appendix (t3b_random.tex, now in the
% supplementary document). The reasons of the manual sets
% follow LMSPS review/manual_metapath_sets.md (revised version).
\begin{widetable}[htbp]
  \centering
  \small
  \caption{Cả 24 tập meta-path: tìm được, thủ công và chọn ngẫu nhiên làm đối chứng. Một meta-path
    được viết từ sản phẩm đi ra
    và đại diện cho cả chính nó lẫn meta-path ngược của nó. Ví dụ, \code{PUP} nối hai sản phẩm được cùng
    một người dùng nhận xét. Tập có dấu * là tập chọn theo xác thực của HAN
    (Mục~\ref{sec:res-random}).}
  \label{tab:sets}
  \begin{tabularx}{\linewidth}{@{}llL@{}}
    \toprule
    Tập & Meta-path & Cách chọn \\
    \midrule
    \multicolumn{3}{@{}l}{\emph{Tìm được: LMSPS, \val{search-runs} lần chạy cho mỗi kích thước
      tập, kết hợp lại}} \\
    \code{a1}, 4 meta-path & \texttt{\val{set-f4a1}} & điểm trung bình tốt nhất, trong một cửa sổ
      \val{search-window-4} tập \\
    \code{a2}, 4 meta-path & \texttt{\val{set-f4a2}} & tập còn lại của cửa sổ đó \\
    \code{c}, 4 meta-path & \texttt{\val{set-f4c}} & tập đồng thuận của bỏ phiếu \\
    \code{a1}, 6 meta-path & \texttt{\val{set-f6a1}} & đơn giản nhất trong \val{search-window-6}
      tập của cửa sổ \\
    \code{a2}, 6 meta-path & \texttt{\val{set-f6a2}} & đơn giản thứ hai \\
    \code{c}, 6 meta-path & \texttt{\val{set-f6c}} & tập đồng thuận của bỏ phiếu \\
    \midrule
    \multicolumn{3}{@{}l}{\emph{Thủ công: thiết kế với sự hỗ trợ của AI theo các nguyên tắc nêu
      rõ, cố định từ trước}} \\
    M1 & \texttt{\val{set-m1}} & một meta-path hành vi, một meta-path nội dung \\
    M2 & \texttt{\val{set-m2}} & mỗi loại bằng chứng một meta-path: mua kèm, đồng nhận xét, danh
      mục \\
    M3 & \texttt{\val{set-m3}} & chỉ hành vi \\
    M4 & \texttt{\val{set-m4}} & chỉ danh mục \\
    M5 & \texttt{\val{set-m5}} & M2 và \code{PPUP}, dùng cho so sánh chính với 4 meta-path \\
    M6 & \texttt{\val{set-m6}} & M5 và hai meta-path rộng hơn, dùng cho so sánh chính với 6 meta-path \\
    M7 & \texttt{\val{set-m7}} & mọi meta-path dài tối đa hai bước, không chọn lọc \\
    M8 & \texttt{\val{set-m8}} & các meta-path dài qua người dùng, như các lần chạy tìm kiếm trước
      đó ưa thích \\
    \midrule
    \multicolumn{3}{@{}l}{\emph{Ngẫu nhiên (đối chứng): chọn đều và không lặp lại từ
      \val{search-cands} ứng viên, cố định từ trước}} \\
    \emph{lần chọn} & \emph{4 meta-path} & \emph{6 meta-path} \\
    1 & \texttt{\val{set-r4d1}} & \texttt{\val{set-r6d1}} \\
    2 & \texttt{\val{set-r4d2}} & \texttt{\val{set-r6d2}}\textsuperscript{*} \\
    3 & \texttt{\val{set-r4d3}} & \texttt{\val{set-r6d3}} \\
    4 & \texttt{\val{set-r4d4}} & \texttt{\val{set-r6d4}} \\
    5 & \texttt{\val{set-r4d5}} & \texttt{\val{set-r6d5}} \\
    \bottomrule
  \end{tabularx}
\end{widetable}


\paragraph{Độ đo.} Mỗi cạnh xác thực và thử nghiệm $(u, p)$ được xếp hạng so với
\val{split-negatives} mẫu âm, lấy mẫu một lần trong số các sản phẩm mà $u$ không có cạnh nối tới.     
Hạng $r$ của $p$ là số ứng viên có điểm cao hơn hoặc bằng điểm của $p$, tính cả chính $p$.
Như vậy, $p$ luôn được xếp sau các ứng viên bằng điểm với nó, nên các trường hợp bằng điểm
được tính bất lợi cho mô hình. Chúng tôi báo cáo Hit@10 ($[r \le 10]$) và NDCG@10
($[r \le 10] / \log_2(r + 1)$) trên \val{split-test} cạnh thử nghiệm. Chúng tôi không báo cáo
Recall@10 và Precision@10 vì với một sản phẩm liên quan cho mỗi lần xếp hạng, chúng lần lượt bằng
Hit@10 và Hit@10$/10$. Hit@10 lấy mẫu là độ đo chúng tôi đã cố
định từ trước, nhưng so với \val{split-negatives} sản phẩm ngẫu nhiên, một sản phẩm phổ biến dễ dàng
thắng: riêng mô hình phổ biến đã đạt \val{pop-hit}. Ngoài ra, các độ đo lấy mẫu có thể xếp thứ tự các mô
hình khác với xếp hạng toàn bộ \citep{krichene2022sampled}. Để kiểm tra, chúng tôi cũng xếp hạng
mỗi cạnh thử nghiệm so với mọi sản phẩm mà người dùng của nó không có cạnh nối tới (gần như toàn bộ
\val{after-products} sản phẩm), rồi báo cáo Hit@10 theo cạnh và Recall@20 theo người dùng, với mọi sản
phẩm thử nghiệm của một người dùng được xếp hạng cùng nhau, trên \val{test-users} người dùng có
cạnh thử nghiệm. Mô hình phổ biến được báo cáo bên cạnh các mô hình trong Bảng~\ref{tab:main}.
Các cạnh thử nghiệm cũng được chia theo số cạnh huấn luyện của người dùng (1--5, 6--20 và từ 21 trở lên).

\paragraph{Thiết lập.} Mỗi mô hình có một cấu hình duy nhất, cố định trước các lần chạy và dùng cho
mọi tập meta-path. Không có gì được tinh chỉnh theo từng tập hay theo kết quả thử nghiệm. Các thiết
lập được chọn để vừa với một GPU \val{env-gpu-mem}\,GB và để huấn luyện ổn định, chứ không được dò
tìm để tối đa độ chính xác. HAN có \val{han-set-num-heads} đầu chú ý \val{han-set-head-dim} chiều
cho mỗi meta-path và embedding người dùng \val{han-set-hidden} chiều. GAT có \val{gat-set-layers}
tầng \val{gat-set-hidden} chiều. metapath2vec học embedding \val{mtv-set-dim} chiều với cửa sổ
\val{mtv-set-window}. Mọi bộ tối ưu đều là Adam, với tốc độ học \val{mtv-set-lr} cho metapath2vec
và \val{han-set-lr} cho tìm kiếm cùng các mô hình còn lại. Việc huấn luyện dừng khi BPR xác thực
không cải thiện quá \val{han-set-min-delta} trong \val{han-set-patience} epoch
(với metapath2vec là \val{mtv-set-patience} lần kiểm tra). Sau đó, điểm lưu tốt nhất được thử
nghiệm một lần. Mọi lần chạy được báo cáo đều dừng sớm: HAN ở epoch
\val{han-stop-min}--\val{han-stop-max} trong tối đa \val{han-set-stage} epoch.

\paragraph{Seed và thống kê.} Phép chia được cố định, nên các seed chỉ thay đổi quá trình huấn
luyện. So sánh chính (\code{a1} so với M5 và M6, trong HAN và trong metapath2vec) và GAT chạy với
seed 1--3, các tập khác với seed~1. Qua ba seed, Hit@10 thử nghiệm của các lần chạy chính dao động
nhiều nhất \val{main-hit-sd-max} (độ lệch chuẩn), nên các chênh lệch nhỏ giữa những tập chạy với
một seed nằm trong nhiễu giữa các seed. Vì vậy, chúng tôi không xếp hạng chi li những tập như vậy. Chúng
tôi so sánh các tập chính bằng kiểm định $t$ hai phía của Welch trên ba seed của mỗi bên
\citep{welch1947generalization}. Với chỉ ba seed mỗi bên, kiểm định này khó phát hiện các chênh lệch nhỏ, nên
chúng tôi cũng nêu mỗi chênh lệch theo đơn vị độ lệch chuẩn giữa các seed. Chỉ Hit@10 lấy mẫu được cố định từ
trước. Các độ đo khác và giá trị $p$ của chúng chỉ mang tính bổ trợ. Các giá trị $p$ này không được hiệu
chỉnh cho việc thực hiện nhiều phép so sánh. Quan hệ giữa các tập dùng hệ số tương quan Spearman với giá
trị $p$ hoán vị từ \val{perm-shuffles} lần xáo trộn. \emph{Tập chọn theo xác thực} là tập có BPR
xác thực thấp nhất.

\paragraph{Phần cứng và chi phí.} Mọi lần chạy đều dùng một GPU NVIDIA \val{env-gpu} với Python
\val{env-python} và PyTorch \val{env-torch} (CUDA \val{env-cuda}). Toàn bộ \val{search-runs-total}
lần chạy tìm kiếm mất khoảng \val{search-total-h}~giờ~GPU (Bảng~\ref{tab:cost}), còn các lần chạy mô
hình mất khoảng \val{models-hours}~giờ~GPU.

% sections/results.tex -- Section 7 (paper_v2_vn, translated in L6 from paper_v2's
% sections/results.tex; terms as in glossary_vn.tex). The details that left it in paper_v2 (the
% paths' ranks, the delta check, the user groups, more paths, the path weights and preferences, HAN
% against metapath2vec) are in the supplementary document (supplement_vn.tex, N31). A few of their
% numbers are given here in the text. Every number is \val{} from numbers.tex. Tables: T4
% (tab:main), T7 (tab:cost); figures: F3 (fig:sets),
% F4 (fig:weights).
\section{Kết quả}
\label{sec:results}

Mọi con số đều là kết quả trên tập thử nghiệm trừ khi có ghi chú khác.

\subsection{Tập tìm được so với tập thủ công}
\label{sec:res-main}

Theo độ đo đã cố định từ trước, \code{a1} thua tập thủ công cùng kích thước trên cả hai mô hình
(Bảng~\ref{tab:main}). Trong HAN nó kém M5 \val{han-4-hit-diffabs} (\val{han-4-hit-relabs}\%,
\val{han-4-hit-pt}) và kém M6 \val{han-6-hit-diffabs} (\val{han-6-hit-relabs}\%,
\val{han-6-hit-pt}), bằng \val{han-4-hit-sdx} và \val{han-6-hit-sdx} lần độ lệch chuẩn giữa các
seed (lấy bên lớn hơn). Trong metapath2vec, chênh lệch rộng hơn: \val{mtv-4-hit-diffabs}
(\val{mtv-4-hit-relabs}\%) và \val{mtv-6-hit-diffabs} (\val{mtv-6-hit-relabs}\%). Khi xếp hạng
toàn bộ, các tập của HAN ngang bằng nhau (\val{han-4-fhit-pt} và \val{han-6-fhit-pt}). Trong
metapath2vec, \code{a1} ngang bằng với $M = 4$ meta-path (\val{mtv-4-fhit-pt}) và thua với $M = 6$ (\val{mtv-6-fhit-pt}). Trong
HAN sự thua kém này cũng đúng với \code{a2} và \code{c} (Hình~\ref{fig:sets}) và ở mọi nhóm người
dùng (chênh lệch từ \val{han-6-hit-d3-diff} đến \val{han-4-hit-d1-diff}).

% t4_main.tex -- T4: the main comparison (paper_v2_vn, translated in L3 from paper_v2's
% tables/t4_main.tex; results.csv via numbers.tex; Welch tests computed by make_numbers.py).
\begin{widetable}[htbp]
  \centering
  \small
  \setlength{\tabcolsep}{4pt}
  \caption{Tập tìm được so với tập thủ công: kết quả trên tập thử nghiệm, trung bình $\pm$ độ lệch
    chuẩn qua 3 seed huấn luyện. $\Delta$: tập tìm được trừ tập thủ công, kèm $p$ của kiểm định
    Welch hai phía. Lấy mẫu: mỗi cạnh thử nghiệm được xếp hạng so với \val{split-negatives} mẫu âm,
    với Hit@10 là độ đo đã cố định từ trước. Toàn bộ: so với mọi sản phẩm mà người dùng không có
    cạnh nối tới, với Recall@20 lấy trung bình theo người dùng.}
  \label{tab:main}
  \begin{tabular*}{\linewidth}{@{\extracolsep{\fill}}llllll@{}}
    \toprule
    & & \multicolumn{2}{c}{Lấy mẫu} & \multicolumn{2}{c}{Xếp hạng toàn bộ} \\
    \cmidrule(lr){3-4}\cmidrule(l){5-6}
    Mô hình & Tập & Hit@10 & NDCG@10 & Hit@10 & Recall@20 \\
    \midrule
    HAN & \code{a1}, 4 meta-path & \val{han-f4a1-hit-mean} $\pm$ \val{han-f4a1-hit-sd}
      & \val{han-f4a1-ndcg-mean} & \val{han-f4a1-fhit-mean} $\pm$ \val{han-f4a1-fhit-sd}
      & \val{han-f4a1-frec-mean} \\
    & M5 & \val{han-m5-hit-mean} $\pm$ \val{han-m5-hit-sd} & \val{han-m5-ndcg-mean}
      & \val{han-m5-fhit-mean} $\pm$ \val{han-m5-fhit-sd} & \val{han-m5-frec-mean} \\
    & $\Delta$ ($p$) & \val{han-4-hit-diff} (\val{han-4-hit-p})
      & \val{han-4-ndcg-diff} (\val{han-4-ndcg-p}) & \val{han-4-fhit-diff} (\val{han-4-fhit-p})
      & \val{han-4-frec-diff} (\val{han-4-frec-p}) \\
    \addlinespace
    & \code{a1}, 6 meta-path & \val{han-f6a1-hit-mean} $\pm$ \val{han-f6a1-hit-sd}
      & \val{han-f6a1-ndcg-mean} & \val{han-f6a1-fhit-mean} $\pm$ \val{han-f6a1-fhit-sd}
      & \val{han-f6a1-frec-mean} \\
    & M6 & \val{han-m6-hit-mean} $\pm$ \val{han-m6-hit-sd} & \val{han-m6-ndcg-mean}
      & \val{han-m6-fhit-mean} $\pm$ \val{han-m6-fhit-sd} & \val{han-m6-frec-mean} \\
    & $\Delta$ ($p$) & \val{han-6-hit-diff} (\val{han-6-hit-p})
      & \val{han-6-ndcg-diff} (\val{han-6-ndcg-p}) & \val{han-6-fhit-diff} (\val{han-6-fhit-p})
      & \val{han-6-frec-diff} (\val{han-6-frec-p}) \\
    \midrule
    metapath2vec & \code{a1}, 4 meta-path & \val{mtv-f4a1-hit-mean} $\pm$ \val{mtv-f4a1-hit-sd}
      & \val{mtv-f4a1-ndcg-mean} & \val{mtv-f4a1-fhit-mean} $\pm$ \val{mtv-f4a1-fhit-sd}
      & \val{mtv-f4a1-frec-mean} \\
    & M5 & \val{mtv-m5-hit-mean} $\pm$ \val{mtv-m5-hit-sd} & \val{mtv-m5-ndcg-mean}
      & \val{mtv-m5-fhit-mean} $\pm$ \val{mtv-m5-fhit-sd} & \val{mtv-m5-frec-mean} \\
    & $\Delta$ ($p$) & \val{mtv-4-hit-diff} (\val{mtv-4-hit-p})
      & \val{mtv-4-ndcg-diff} (\val{mtv-4-ndcg-p}) & \val{mtv-4-fhit-diff} (\val{mtv-4-fhit-p})
      & \val{mtv-4-frec-diff} (\val{mtv-4-frec-p}) \\
    \addlinespace
    & \code{a1}, 6 meta-path & \val{mtv-f6a1-hit-mean} $\pm$ \val{mtv-f6a1-hit-sd}
      & \val{mtv-f6a1-ndcg-mean} & \val{mtv-f6a1-fhit-mean} $\pm$ \val{mtv-f6a1-fhit-sd}
      & \val{mtv-f6a1-frec-mean} \\
    & M6 & \val{mtv-m6-hit-mean} $\pm$ \val{mtv-m6-hit-sd} & \val{mtv-m6-ndcg-mean}
      & \val{mtv-m6-fhit-mean} $\pm$ \val{mtv-m6-fhit-sd} & \val{mtv-m6-frec-mean} \\
    & $\Delta$ ($p$) & \val{mtv-6-hit-diff} (\val{mtv-6-hit-p})
      & \val{mtv-6-ndcg-diff} (\val{mtv-6-ndcg-p}) & \val{mtv-6-fhit-diff} (\val{mtv-6-fhit-p})
      & \val{mtv-6-frec-diff} (\val{mtv-6-frec-p}) \\
    \midrule
    Phổ biến & -- & \val{pop-hit} & \val{pop-ndcg} & \val{pop-fhit} & \val{pop-frec} \\
    \bottomrule
  \end{tabular*}
\end{widetable}


\subsection{Đối chứng ngẫu nhiên}
\label{sec:res-random}

Để biết khoảng cách giữa tìm kiếm và thiết kế thủ công lớn đến đâu so với ảnh hưởng của việc chọn
meta-path, Hình~\ref{fig:sets} đặt các tập tìm được và các tập thủ công giữa cả 24 tập, cùng với các tập
chọn ngẫu nhiên làm đối chứng. Trong tám tập của mỗi kích thước được lấy
từ cùng \val{search-cands} ứng viên, gồm ba tập tìm được và năm tập ngẫu nhiên, HAN xếp các tập
tìm được ở vị trí \val{han-4-pos-hit} với $M = 4$ meta-path và \val{han-6-pos-hit} với $M = 6$, giống như
BPR xác thực. Khi xếp hạng toàn bộ, các tập này đứng ở vị trí \val{han-4-pos-fhit} với $M = 4$ và
\val{han-6-pos-fhit} với $M = 6$. Các khoảng cách này nằm trong nhiễu giữa các seed: tập tìm được tốt
nhất đạt \val{han-f4c-hit-s1} với $M = 4$ và \val{han-f6a1-hit-s1} với $M = 6$, so với
\val{han-r4-hit-min} và \val{han-r6-hit-min} của những tập ngẫu nhiên yếu nhất. Vì vậy, các vị trí này cho
thấy tìm kiếm không vượt được chính các ứng viên của nó, chứ không cho thấy mọi tập ngẫu nhiên đều vượt nó.
Các tập thủ công đạt kết quả xấp xỉ một tập ngẫu nhiên trung bình: M5 đạt \val{han-m5-hit-mean} so với
trung bình \val{han-r4-hit-mean} của các tập ngẫu nhiên $M = 4$ meta-path (khoảng
\val{han-r4-hit-min}--\val{han-r4-hit-max}) và M6 đạt \val{han-m6-hit-mean} so với
\val{han-r6-hit-mean} (\val{han-r6-hit-min}--\val{han-r6-hit-max}).

% f3_sets.tex -- F3: HAN test Hit@10 of all 24 sets by family (paper_v2_vn, translated in L3 from
% paper_v2's figures/f3_sets.tex). Data: f3_*.dat, written by make_numbers.py from results.csv
% (seed 1 of every set).
\begin{widefigure}[htbp]
  \centering
  \begin{tikzpicture}
    \begin{axis}[
        width=0.9\linewidth, height=5.4cm,
        xmin=0.45, xmax=3.6, ymin=0.612, ymax=0.712,
        xtick={1,2,3}, xticklabels={2--3 meta-path, 4 meta-path, 6 meta-path},
        ylabel={Hit@10 thử nghiệm của HAN (lấy mẫu)},
        yticklabel style={/pgf/number format/fixed, /pgf/number format/precision=2},
        ytick={0.62,0.64,0.66,0.68,0.70},
        tick label style={font=\scriptsize}, label style={font=\footnotesize},
        ymajorgrids, grid style={black!8}, axis line style={black!60}, tick style={black!60},
        legend style={font=\scriptsize, at={(0.02,0.98)}, anchor=north west, draw=none,
          fill=white, fill opacity=0.85, text opacity=1},
        legend cell align=left]
      \draw[black!20] (axis cs:1.5,0.612) -- (axis cs:1.5,0.712);
      \draw[black!20] (axis cs:2.5,0.612) -- (axis cs:2.5,0.712);
      \addplot[only marks, mark=*, mark size=2.6pt, mark options={fill=cSearch, draw=white,
          line width=0.4pt}] table[x=x, y=hit] {figures/f3_searched.dat};
      \addlegendentry{tìm được (LMSPS)}
      \addplot[only marks, mark=square*, mark size=2.4pt, mark options={fill=cManual, draw=white,
          line width=0.4pt}] table[x=x, y=hit] {figures/f3_manual.dat};
      \addlegendentry{thủ công}
      \addplot[only marks, mark=triangle*, mark size=3pt, mark options={fill=cRandom, draw=white,
          line width=0.4pt}] table[x=x, y=hit] {figures/f3_random.dat};
      \addlegendentry{chọn ngẫu nhiên}
      \addplot[only marks, mark=none, point meta=explicit symbolic,
          visualization depends on={value \thisrow{anchor}\as\labelanchor},
          nodes near coords, nodes near coords style={font=\scriptsize, text=black!80,
            anchor=\labelanchor, inner sep=2.5pt}]
        table[x=x, y=y, meta=label] {figures/f3_labels.dat};
    \end{axis}
  \end{tikzpicture}
  \Description{Biểu đồ điểm Hit@10 thử nghiệm của HAN cho cả 24 tập meta-path, nhóm theo kích thước tập, mỗi loại tập một hình dạng và màu điểm riêng. Các tập tìm được nằm ở phía thấp của các tập lấy từ các ứng viên. Các tập thủ công nằm giữa các tập chọn ngẫu nhiên.}
  \caption{Hit@10 thử nghiệm của HAN cho cả 24 tập meta-path, mỗi tập một seed (độ lệch chuẩn qua
    ba seed của các tập chính nhiều nhất là \val{han-m6-hit-sd}). Ở seed này, các tập tìm được đứng
    cuối trong số các tập lấy từ cùng \val{search-cands} ứng viên ở cả hai kích thước, dù với chênh
    lệch nằm trong nhiễu giữa các seed. M5 và M6 nằm gần giữa các tập chọn ngẫu nhiên. Để so sánh:
    mô hình phổ biến \val{pop-hit}, GAT \val{gat-hit-mean}.}
  \label{fig:sets}
\end{widefigure}


Do đó, việc chọn meta-path ít ảnh hưởng đến HAN: cả \val{han-4-nsets} tập $M = 4$ meta-path chỉ chênh
nhau trong phạm vi \val{han-4-span} Hit@10, còn cả \val{han-6-nsets} tập $M = 6$ meta-path chênh nhau trong phạm vi
\val{han-6-span}. Chọn theo xác thực là hiệu quả: trên 24 tập, BPR xác thực và Hit@10 thử nghiệm có
hệ số tương quan Spearman \val{han-valtest-rho} (\val{han-valfull-rho} với Hit@10 xếp hạng toàn
bộ). Tập chọn theo xác thực vốn là một tập ngẫu nhiên đạt \val{han-r6d2-hit-s1}, mức mà không tập
nào khác vượt được theo độ đo đã cố định từ trước. Như vậy, lựa chọn rẻ nhất là chọn ngẫu nhiên vài tập
rồi giữ tập có điểm xác thực tốt nhất. Trong metapath2vec, so sánh với ngẫu nhiên cho kết quả trái
chiều: các tập tìm được đứng \val{mtv-4-pos-hit} trong tám tập với $M = 4$ meta-path, trên tập ngẫu nhiên
trung bình, nhưng \val{mtv-6-pos-hit} với $M = 6$ meta-path, dưới mức đó. Hai mô hình cũng ưa thích các
tập khác nhau: trên 23 tập có meta-path qua người dùng, giá trị Hit@10 của chúng có hệ số tương
quan Spearman chỉ \val{han-mtv-rho}.

\subsection{Chi phí so với hiệu quả}
\label{sec:res-cost}

Thiết kế thủ công không dùng GPU và thời gian của nó cũng không được ghi lại (Bảng~\ref{tab:cost}).
Các công trình trước cũng chỉ mô tả chi phí này bằng lời \citep[][Phụ lục~D]{han2020genetic,li2024longrange}.
LMSPS còn bỏ thời gian thiết kế thủ công ra khỏi phép so sánh hiệu quả của mình, cùng với thời gian
tìm kiếm của chính nó \citep[Phụ lục~E.1]{li2024longrange}. Tuy vậy, chi phí của tìm kiếm cho chúng
tôi một mốc để so sánh: toàn bộ \val{search-runs-total} lần chạy tìm kiếm mất khoảng
\val{search-total-h}~giờ~GPU, chưa kể thời gian điều chỉnh mã nguồn. Thiết kế thủ công rẻ hơn tìm
kiếm chừng nào việc thiết kế mất ít hơn khoảng thời gian này. Chúng tôi cho rằng điều đó hợp lý trên
một lược đồ chỉ có ba kiểu nút, nơi các tập thủ công được thiết kế theo bốn nguyên tắc trong một
buổi, cùng các phép kiểm tra độ dày đặc chỉ mất vài phút CPU. Trong khi đó, tìm kiếm thu được những
tập kém các tập thủ công \val{han-4-hit-relabs}\% và \val{han-6-hit-relabs}\% theo độ đo đã cố định
từ trước và chỉ ngang bằng chúng khi xếp hạng toàn bộ. Như vậy, trên lược đồ này, khoản chi phí tìm
kiếm chưa đem lại cho HAN kết quả vượt trội hơn thiết kế thủ công. Đối chứng ngẫu nhiên vốn không
tốn chi phí, cũng đạt kết quả ít nhất ngang bằng cả hai cách. Khi xếp hạng toàn bộ, trung bình kết
quả của các tập ngẫu nhiên ít nhất cũng cao bằng, dù phân tán rộng hơn
(\val{han-r4-fhit-min}--\val{han-r4-fhit-max} và \val{han-r6-fhit-min}--\val{han-r6-fhit-max}).

% t7_cost.tex -- T7: cost and performance of the three ways to choose HAN's paths (paper_v2_vn,
% translated in L3 from paper_v2's tables/t7_cost.tex). Search cost from the 60 search logs (150
% epochs at the logged epoch time plus the feature propagation); results from results.csv; all via
% numbers.tex. The cost of the hand design is stated in words: none of the reference papers
% measures it.
\begin{widetable}[htbp]
  \centering
  \small
  \caption{Chi phí và kết quả của từng cách chọn meta-path cho HAN: Hit@10 trên tập thử
    nghiệm, lấy mẫu và xếp hạng toàn bộ. Thủ công và bằng tìm kiếm: trung bình qua 3 seed. Ngẫu
    nhiên: trung bình của 5 lần chọn, mỗi lần một seed. Chênh lệch tương đối được tính so với tập
    thủ công cùng kích thước. Huấn luyện HAN trên tập đã chọn tốn như nhau cho cả ba cách, trung vị
    \val{han-run-min}~phút~GPU mỗi lần chạy.}
  \label{tab:cost}
  \begin{tabularx}{\linewidth}{@{}lLlll@{}}
    \toprule
    & & \multicolumn{2}{l}{Hit@10 lấy mẫu} & Hit@10 toàn bộ \\
    \cmidrule(lr){3-4}\cmidrule(l){5-5}
    Cách chọn & Chi phí lựa chọn & 4 / 6 meta-path & so với thủ công & 4 / 6 meta-path \\
    \midrule
    Thủ công (M5, M6) & không dùng GPU: một buổi thiết kế có AI hỗ trợ (không ghi lại thời gian)
      cùng các phép kiểm tra độ dày đặc trên một đồ thị nhỏ, mất vài phút CPU
      & \val{han-m5-hit-mean} / \val{han-m6-hit-mean} & --
      & \val{han-m5-fhit-mean} / \val{han-m6-fhit-mean} \\
    Bằng tìm kiếm (\code{a1}) & \val{search-runs} lần chạy cho mỗi kích thước, trung vị
      \val{search-run-min-4} / \val{search-run-min-6}~phút~GPU mỗi lần:
      \val{search-total-h}~giờ~GPU & \val{han-f4a1-hit-mean} / \val{han-f6a1-hit-mean}
      & \val{han-4-hit-rel}\% / \val{han-6-hit-rel}\%
      & \val{han-f4a1-fhit-mean} / \val{han-f6a1-fhit-mean} \\
    Ngẫu nhiên (5 lần chọn) & không tốn & \val{han-r4-hit-mean} / \val{han-r6-hit-mean}
      & \val{han-r4-vs-manual-rel}\% / \val{han-r6-vs-manual-rel}\%
      & \val{han-r4-fhit-mean} / \val{han-r6-fhit-mean} \\
    \bottomrule
  \end{tabularx}
\end{widetable}


\subsection{HAN làm gì với các meta-path}
\label{sec:res-weights}

Hình~\ref{fig:weights} so sánh mức ưa thích của tìm kiếm đối với từng meta-path với trọng số mà
HAN gán cho nó. Hai bên thống nhất về các meta-path qua người dùng. \code{PUPUP} vốn nối các sản
phẩm qua hai người dùng, có mặt trong cả \val{searched-with-pupup} tập tìm được và \code{PUPPUP}
trong \val{searched-with-puppup} tập. Cả hai đều xếp gần đầu trong các ứng viên của tìm kiếm.
HAN không bao giờ loại chúng hay \code{PUP}, đồng thời gán cho cả ba trọng số cao hơn mức chia
đều. Hai bên bất đồng về danh mục. Mỗi tập tìm được chứa \val{searched-empty-min} đến
\val{searched-empty-max} meta-path danh mục gần rỗng mà tìm kiếm xếp hạng cao, nhưng HAN chỉ giữ
một trong số đó cho mỗi tập với trọng số \val{han-empty-kept-min} đến \val{han-empty-kept-max}
(\val{han-empty-kept-f4c} với \code{c} ở $M = 4$ meta-path), nên các tập tìm được $M = 6$ vận hành
như các tập $M = 4$. Các meta-path danh mục dày đặc của các tập thủ công xếp hạng thấp trong
tìm kiếm nhưng nhận được khoảng mức chia đều hoặc hơn trong HAN (trung vị \val{kind-cat-weight}).
Không tập tìm được nào chứa một meta-path danh mục đến được phần lớn sản phẩm.

% f4_weights.tex -- F4: HAN's path weights against the search's preference, by kind of path
% (paper_v2_vn, translated in L3 from paper_v2's figures/f4_weights.tex). Data: f4_*.dat, written
% by make_numbers.py (search records; HAN seed-1 logs).
\begin{figure}[htbp]
  \centering
  \begin{tikzpicture}
    \begin{axis}[
        width=\linewidth, height=5.6cm,
        xmin=0, xmax=33, ymin=-0.1, ymax=2.05, xtick={1,8,16,24,32},
        xlabel={Mức ưa thích của tìm kiếm: hạng $\alpha$ trung bình (1 = tốt nhất trong
          \val{search-cands})},
        ylabel={Trọng số của HAN (1 = mức chia đều)},
        tick label style={font=\scriptsize}, label style={font=\footnotesize},
        ymajorgrids, grid style={black!8}, axis line style={black!60}, tick style={black!60},
        legend style={font=\scriptsize, at={(0.5,1.02)}, anchor=south, draw=none,
          legend columns=2, /tikz/every even column/.append style={column sep=6pt}},
        legend cell align=left]
      \draw[black!45, dashed] (axis cs:0,1) -- (axis cs:33,1);
      \addplot[only marks, mark=*, mark size=2.6pt, mark options={fill=cSearch, draw=white,
          line width=0.4pt}] table[x=rank, y=weight] {figures/f4_users.dat};
      \addlegendentry{qua người dùng (\val{kind-users-paths} meta-path)}
      \addplot[only marks, mark=square*, mark size=2.4pt, mark options={fill=cManual, draw=white,
          line width=0.4pt}] table[x=rank, y=weight] {figures/f4_empty.dat};
      \addlegendentry{danh mục gần rỗng (\val{kind-empty-paths})}
      \addplot[only marks, mark=triangle*, mark size=3pt, mark options={fill=cRandom, draw=white,
          line width=0.4pt}] table[x=rank, y=weight] {figures/f4_cat.dat};
      \addlegendentry{danh mục khác (\val{kind-cat-paths})}
      \addplot[only marks, mark=diamond*, mark size=3pt, mark options={fill=cOther, draw=white,
          line width=0.4pt}] table[x=rank, y=weight] {figures/f4_pp.dat};
      \addlegendentry{chỉ sản phẩm (\val{kind-pp-paths})}
      \addplot[only marks, mark=none, point meta=explicit symbolic,
          visualization depends on={value \thisrow{anchor}\as\labelanchor},
          nodes near coords, nodes near coords style={font=\scriptsize\ttfamily, text=black!80,
            anchor=\labelanchor, inner sep=2.5pt}]
        table[x=x, y=y, meta=label] {figures/f4_labels.dat};
    \end{axis}
  \end{tikzpicture}
  \Description{Biểu đồ phân tán, mỗi điểm là một meta-path: mức ưa thích trung bình của tìm kiếm đối với meta-path so với trọng số trung bình mà HAN gán cho nó. Cả hai đều ưa các meta-path qua người dùng. Các meta-path danh mục gần rỗng được tìm kiếm ưa thích nhưng gần như không có trọng số trong HAN. Các meta-path danh mục dày đặc bị tìm kiếm xếp hạng thấp nhưng được HAN sử dụng.}
  \caption{HAN làm gì với các meta-path: mỗi điểm là một meta-path, đặt theo mức ưa thích trung
    bình của tìm kiếm đối với nó (cả hai kích thước tập) và trọng số trung bình HAN gán cho nó trên
    các tập chứa nó (seed 1). Hai bên thống nhất về các meta-path qua người dùng. Hai bên bất đồng
    về các meta-path danh mục gần rỗng, vốn được tìm kiếm xếp hạng cao còn HAN phần lớn loại bỏ
    (\val{kind-empty-dropped} trong \val{kind-empty-uses} lần dùng có trọng số dưới
    \val{weight-drop}) và về các meta-path danh mục dày đặc, vốn bị tìm kiếm xếp hạng thấp còn HAN
    sử dụng.}
  \label{fig:weights}
\end{figure}


Trong các tập $M = 4$ và $M = 6$ meta-path, sự thiếu vắng này phân biệt các tập tìm được với tất cả các
tập khác. Ở cả hai kích thước, mọi tập có một meta-path danh mục đến được gần như mọi sản phẩm
đều vượt mọi tập không có meta-path như vậy: \val{han-4-dense-min}--\val{han-4-dense-max} so với
\val{han-4-nodense-min}--\val{han-4-nodense-max} với $M = 4$ (cặp gần nhau nhất nằm trong
nhiễu giữa các seed) và \val{han-6-dense-min}--\val{han-6-dense-max} so với \val{han-6-nodense-max}
với $M = 6$. Trên \val{pref-cand-n} tập lấy từ các ứng viên, tìm kiếm càng ưa thích các meta-path của
một tập thì HAN làm càng kém với tập đó (hệ số Spearman \val{pref-cand-rho} giữa Hit@10 thử nghiệm và
hạng $\alpha$ trung bình của các meta-path, trong đó hạng càng thấp nghĩa là càng được ưa thích,
\val{pref-cand-pt}). Quan hệ này dựa trên các tập tìm được: nếu chỉ xét mười tập chọn ngẫu nhiên,
hệ số này chỉ còn \val{pref-random-rho} (\val{pref-random-pt}).

\paragraph{Nơi tập trung sự thua kém.} Một bước đơn lẻ trên cây phân cấp danh mục, \code{PCCP} nối
\val{bg-pccp-reach-products-pct}\% số sản phẩm với một sản phẩm khác, phần lớn là DVD và video vốn
chiếm \val{bg-pccp-reach-test-pct}\% số cạnh thử nghiệm. Trên các sản phẩm mà \code{PCCP} không đến
được, các tập tìm được kém các tập có meta-path danh mục dày đặc \val{bg-4-gapabs-no} Hit@10 với $M = 4$ meta-path
và \val{bg-6-gapabs-no} với $M = 6$ meta-path (seed~1). Trên các sản phẩm nó đến được, các tập tìm được chỉ kém
\val{bg-4-gapabs-reach} và \val{bg-6-gapabs-reach}. Như vậy, sự thua kém chủ yếu nằm ở sách và
nhạc, những sản phẩm mà các meta-path danh mục của các tập tìm được là rỗng.

\paragraph{Vì sao các tập tìm được thua.} Phép kiểm tra trên chỉ ra các meta-path danh mục gần
rỗng. Trong tìm kiếm, một sản phẩm mà meta-path như vậy không đến được nhận một đặc trưng toàn số 0
cho meta-path đó. Số 0 ấy vẫn mang thông tin: các sản phẩm mà một bước trên cây phân cấp đến được,
phần lớn là DVD và video, chiếm \val{pccp-reach-train-pct}\% số cạnh huấn luyện nhưng chỉ
\val{bg-pccp-reach-products-pct}\% số sản phẩm và do đó cũng chừng ấy phần các mẫu âm ngẫu nhiên,
nên các meta-path rỗng làm lộ nhóm sản phẩm và một phần độ phổ biến của sản phẩm. Trong HAN, tín
hiệu này phần lớn biến mất: mỗi sản phẩm giữ đặc trưng riêng của nó qua một cạnh tự nối trong mọi
đồ thị, nên một kênh rỗng chỉ lặp lại các đặc trưng đó. HAN loại bỏ phần lớn các kênh này. Thứ mà
các tập tìm được khi đó còn thiếu, một meta-path danh mục đến được sách và nhạc, chính là chỗ chúng
thua nhiều nhất. Tìm kiếm cũng có thể tận dụng một meta-path dày đặc kém hơn HAN, vì nó lấy trung
bình đặc trưng của meta-path trên mọi điểm cuối, trong khi HAN chỉ chú ý trên \val{han-set-topk} láng
giềng được chọn. Tuy nhiên, phép kiểm tra theo nhóm sản phẩm chỉ xác định được nơi xảy ra sự thua
kém chứ chưa cho thấy một meta-path danh mục dày đặc sẽ loại bỏ được nó.

\subsection{Các mô hình khác}
\label{sec:res-models}

GAT vượt cả hai mô hình dùng meta-path trên mọi độ đo và ở mọi nhóm người dùng. Nó đạt Hit@10
\val{gat-hit-mean}, so với
\val{han-r6d2-hit-s1} của tập chọn theo xác thực của HAN và \val{mtv-m5-hit-mean} của metapath2vec.
Khi xếp hạng toàn bộ, nó đạt \val{gat-fhit-mean} so với \val{han-r6d2-fhit-s1} của HAN. Mức dẫn trước
của nó so với HAN lớn nhất với những người dùng có ít cạnh huấn luyện: \val{gat-lead-d1} với 1--5
cạnh, \val{gat-lead-d2} với 6--20 và \val{gat-lead-d3} từ 21 cạnh trở lên. Khác với hai
mô hình kia, GAT dựng mỗi người dùng từ các sản phẩm mà người đó đã nhận xét, trong khi HAN và
metapath2vec học một embedding tự do cho mỗi người dùng, embedding này chỉ nhận được ít tín hiệu từ
một vài cạnh. Điều này phù hợp với phát hiện của HGB rằng một GAT đồng nhất với đầu vào phù hợp là
một mô hình cơ sở mạnh \citep{lv2021hgb}. Điều này cũng gợi ý rằng ở đây cách biểu diễn người dùng
quan trọng hơn việc chọn những meta-path phía sản phẩm nào. Các meta-path phía người dùng vốn nằm
ngoài thiết kế của chúng tôi sẽ kiểm chứng được điều này.

HAN vượt metapath2vec về Hit@10 lấy mẫu trên cả
\val{han-wins-sampled} tập nhưng metapath2vec thắng khi xếp hạng toàn bộ ở \val{mtv-wins-full} tập,
nên thứ tự của hai mô hình phụ thuộc vào độ đo. Mọi mô hình đều vượt mức
sàn của mô hình phổ biến (\val{pop-hit} lấy mẫu, \val{pop-fhit} khi xếp hạng toàn bộ), trừ
metapath2vec trên M4 vốn đạt \val{mtv-m4-hit-s1}: tập này không có meta-path nào qua người
dùng nên người dùng không được huấn luyện.

\subsection{Độ nhạy}
\label{sec:res-sensitivity}

Với giai đoạn khởi động trọng số meta-path \val{wu5-warmup} epoch thay vì
\val{han-set-beta-warmup} (seed~1, cả 24 tập), HAN huấn luyện kém hơn về trung bình
(\val{wu5-diff-mean} Hit@10, từ \val{wu5-diff-min} đến \val{wu5-diff-max}), loại nhiều meta-path
hơn (\val{wu5-dropped} trọng số dưới \val{weight-drop} so với \val{han-dropped}) và sắp thứ tự các
tập khác đi (Spearman \val{wu5-rho} so với các lần chạy chính). Khi đó, các tập tìm được đứng
\val{wu5-4-pos} trong tám tập ở cả hai kích thước thay vì đứng cuối. Với $M = 6$ meta-path, phép so
sánh ở seed 1 đảo ngược: \code{a1} đạt \val{wu5-f6a1-hit} so với \val{wu5-m6-hit} của M6 vốn chỉ
giữ \val{wu5-m6-kept} trong $M = 6$ meta-path của nó. Phép kiểm tra này dùng một seed và không được lên
kế hoạch từ trước nên nó không thay đổi kết quả chính nhưng cho thấy phép so sánh phụ thuộc vào
cách HAN học trọng số meta-path: các kết luận của chúng tôi đúng với giai đoạn khởi động
\val{han-set-beta-warmup}~epoch.

% sections/discussion.tex -- the conclusion (paper_v2_vn). Since 20 November (N26) there is no
% Discussion or Threats section: their main points are in Section 7 (why the searched sets lose, why
% GAT leads) and in the conclusion below (the cheapest choice, who should judge a search, the
% limitations). The two minor limitations, in an appendix from N26, are back here since N31. Every
% number is \val{} from numbers.tex.
\section{Kết luận}
\label{sec:conclusion}

Trên một phiên bản Amazon-Meta đã làm sạch với ba kiểu nút và \val{search-cands} meta-path ứng
viên, kết quả của tìm kiếm meta-path tự động khá gần với thiết kế thủ công cho gợi ý sản phẩm dựa
trên HAN: các tập mà một phép tìm kiếm gồm \val{search-runs-total} lần chạy chọn ra kém các tập thủ
công \val{han-6-hit-relabs}--\val{han-4-hit-relabs}\% theo độ đo đã cố định từ trước và ngang bằng
chúng khi xếp hạng toàn bộ. Sự thua kém chủ yếu nằm ở những sản phẩm mà các meta-path danh mục của
các tập tìm được không đến được. Vì vậy, nên chọn cách nào tùy thuộc vào bối cảnh. Nếu có chuyên
gia để thiết kế được meta-path trong thời gian ngắn hơn thời gian tìm kiếm, ở đây sẽ khoảng
\val{search-total-h}~giờ~GPU, thiết kế thủ công là lựa chọn tốt hơn. Nếu không có chuyên gia,
tìm kiếm tự động cho các tập gần với thiết kế thủ công mà không cần kiến thức chuyên môn. Tuy
nhiên, đối chứng ngẫu nhiên cho thấy trên lược đồ nhỏ này việc chọn meta-path ít ảnh hưởng: các
tập chọn ngẫu nhiên từ cùng các ứng viên vốn không tốn chi phí, đạt kết quả ít nhất ngang bằng
cả hai cách (Mục~\ref{sec:res-random}). Lợi thế của tìm kiếm (nếu có) nhiều khả năng nằm ở những
lược đồ lớn hơn, nơi cả thiết kế thủ công lẫn chọn ngẫu nhiên đều khó, điều mà nghiên cứu này chưa
kiểm chứng. Một GAT không phân biệt kiểu vượt cả hai mô hình dùng meta-path.

Bài học chính là một phương pháp tìm kiếm cho một mô hình gợi ý nên được đánh giá bằng chính mô hình
gợi ý đó, như RMS đã làm \citep{ning2022automatic}, hoặc ít nhất được kiểm tra so với các tập chọn
ngẫu nhiên từ chính các ứng viên của nó. Các tập mà mạng riêng của LMSPS đánh giá là tốt nhất không
phải là các tập tốt nhất đối với HAN. Mức cải thiện của LMSPS trên HAN
\citep[Bảng~4]{li2024longrange} được đo trên bài toán phân loại nút, nơi mỗi nút được mã hóa từ vùng
lân cận của nó. Trong bài toán dự đoán liên kết của chúng tôi, phần lớn tham số của HAN là các
embedding người dùng tự do mà không meta-path nào định hình.

\paragraph{Hạn chế.} Nghiên cứu dùng một bộ dữ liệu, một bài toán và một lược đồ nhỏ, chỉ với các
meta-path phía sản phẩm. Phép chia ngẫu nhiên duy nhất bỏ qua ngày nhận xét, điều cho phép mô
hình học từ các tương tác về sau và có thể làm thay đổi thứ hạng của các mô hình
\citep{ji2023leakage,meng2020splitting}. Mọi tập dùng chung một cấu hình không tinh chỉnh. Khi giai
đoạn khởi động trọng số meta-path ngắn hơn, tập tìm được lại vượt tập thủ công ở sáu meta-path,
trong một phép kiểm tra chỉ dùng một seed (Mục~\ref{sec:res-sensitivity}). Phép tìm kiếm của chúng
tôi nhỏ hơn so với bài báo gốc vì được định cỡ cho một GPU. Các lần chạy của nó đồng thuận hơn mức
ngẫu nhiên, nhưng một phép tìm kiếm lớn hơn có thể chọn khác. Các tập thủ công được thiết kế theo các nguyên
tắc nêu rõ chứ không phải bởi một chuyên gia về lĩnh vực này, nên một chuyên gia có thể chọn những
tập khác. Việc huấn luyện HAN chưa hội tụ hẳn: BPR xác thực còn cải thiện tới \val{han-gain-max}
trong mười epoch cuối, nhưng với các tập tìm được thì nhiều nhất chỉ \val{searched-gain-max}. Do
đó, huấn luyện lâu hơn nhiều khả năng sẽ nới rộng hơn là thu hẹp khoảng cách giữa hai cách. Sản
phẩm chỉ liên kết với các danh mục lá (Mục~\ref{sec:data}). Nếu liên kết mọi danh mục cấp trên, các
ứng viên và lựa chọn của tìm kiếm sẽ khác.

Hướng nghiên cứu tiếp theo gồm: (1) liệu một meta-path danh mục dày đặc có thu hẹp
được khoảng cách không, (2) một phép tìm kiếm được đánh giá bằng chính mô hình gợi ý của nó, như
RMS, (3) các meta-path phía người dùng, (4) truy vấn top-$k$ của \citet{zhu2020topk}, (5) các lược
đồ lớn hơn, các bộ dữ liệu khác và các phép chia theo ngày, nơi tìm kiếm có thể đem lại lợi ích.


% sections/statements.tex -- the code and data statement and the AI statement (paper_v2_vn,
% translated in L7 from paper_v2's sections/statements.tex. The AI statement is the author's own
% wording, N28.) The details are in the Reproducibility section of supplement_vn.tex (N31). For a
% repository the code paragraph says "bộ mã nguồn", as the AI statement does (the author's choice,
% not the glossary's word).
\section*{Mã nguồn và dữ liệu}
Mã nguồn được công bố công khai: \url{https://github.com/phananh85327/LMSPS} chứa quy trình xử lý
dữ liệu và tìm kiếm, còn \url{https://github.com/phananh85327/han-improve} chứa các mô hình và phần
đánh giá. Các kết quả đến từ commit \texttt{\val{repo-lmsps-commit}} của bộ mã nguồn thứ nhất và
commit \texttt{\val{repo-han-commit}} của bộ mã nguồn thứ hai. Bản ghi của mọi lần chạy được lưu cùng
hai bộ mã nguồn này. Trong bộ thứ nhất, script \code{hgb/generate_data.bat} dựng lại dữ liệu đã làm
sạch từ tệp thô của SNAP (vốn chúng tôi không phân phối lại) và script \code{paper/make_numbers.py}
tính mọi con số, bảng và hình của bài báo từ các bản ghi đó.

\section*{Sử dụng trợ lý AI}
Chúng tôi đã sử dụng một trợ lý AI, Claude (Anthropic). Các tập meta-path thủ công được thiết kế với
sự hỗ trợ của nó theo các nguyên tắc ở Mục~\ref{sec:setup}, trước khi có bất kỳ kết quả nào. Nó đã
rà soát mã nguồn, đề xuất hướng đi và sửa lỗi. Nó cũng hỗ trợ viết script tính các con
số, bảng và hình của bài báo từ các bản ghi lần chạy. Tác giả đã đề xuất nghiên cứu,
cùng các phương pháp nền tảng, bộ dữ liệu và các bài báo tham chiếu, trong một đề cương dự án
dự định so sánh các meta-path được tìm ra với các meta-path chọn thủ công trên HAN và metapath2vec,
với GAT làm đối chứng. Tác giả cũng đã viết các phiên bản đầu tiên của cả hai bộ mã nguồn,
điều chỉnh mã gốc của LMSPS và HAN cho bài toán dự đoán liên kết. Mọi thực nghiệm đều do tác giả chạy.
Tác giả đã rà soát toàn bộ công việc này và chịu trách nhiệm hoàn toàn về nội dung.


\begin{acks}
Tác giả xin cảm ơn TS.~Ngô Dương Hà và TS.~Phạm Nguyễn Huy Phương đã hướng dẫn. Nghiên cứu không
nhận được tài trợ nào.
\end{acks}

% The reference list stays in English, as the works were published (N7).
\begin{otherlanguage*}{english}
\bibliographystyle{ACM-Reference-Format}
\bibliography{references,notes_extra,reading_list}
\end{otherlanguage*}

\end{document}
